\documentclass[12pt]{amsart}
\usepackage{TheBigPreamble}

\usepackage{latexsym}
\usepackage{amssymb}
\usepackage{amsmath}
\usepackage{mathrsfs}

\theoremstyle{plain} % for future references
\theorembodyfont{\sl}

\newtheorem{theorem}[definition]{Theorem}
\newtheorem{lemma}[definition]{Lemma}
\newtheorem{corollary}[definition]{Corollary}
\newtheorem{proposition}[definition]{Proposition}
\newtheorem{amplification}[definition]{Amplification}
\newtheorem{criterion}[definition]{Criterion}
\newtheorem{algorithm}[definition]{Algorithm}
\newtheorem{claim}[definition]{Claim} %\renewcommand{\theclaim}{}
%\newtheorem{conjecture}[definition]{Conjecture}

\newcommand{\Ker}{\mathrm{ker}}
\newcommand{\End}{\mathrm{End}}
\renewcommand{\Re}{\mathrm{Re }}
\newcommand{\im}{\mathrm{im }}
\newcommand{\ann}{\mathrm{Ann}}
\newcommand{\Ann}{\mathrm{Ann}}
\newcommand{\ad}{\mathrm{ad}}
\renewcommand{\hom}{\mathrm{Hom}}
\newcommand{\Hom}{\mathrm{Hom}}
\newcommand{\rk}{\mathrm{rk }}
\newcommand{\ch}{\mathrm{ch}}
\newcommand{\res}{\mathrm{Res}}
\newcommand{\Ob}{\mathrm{Ob}}
\newcommand{\ind}{\mathrm{Ind}}
\newcommand{\Ind}{\mathrm{Ind}}
\newcommand{\GL}{{GL}}
\newcommand{\id}{\mathrm{id}}
\renewcommand{\sl}{\mathrm{sl}}
\newcommand{\gl}{\mathrm{gl}}
\renewcommand{\sp}{\mathrm{sp}}
\newcommand{\so}{\mathrm{so}}
\newcommand{\sign}{\mathrm{sign}}
\renewcommand{\span}{\mathrm{span}}

\newcommand {\ZZ} {\mathbb {Z}}
\newcommand {\CC} {\mathbb {C}}
\newcommand {\EE} {\mathbf {E}}
\newcommand {\RR} {\mathbb {R}}
\newcommand {\PP} {\mathbb {P}}
\newcommand {\HH} {\mathcal {H}}
\newcommand {\NN} {\mathbb {N}}
\renewcommand{\ss}{\mathfrak{s}}
\newcommand{\kk}{\mathfrak{k}}
\renewcommand{\aa}{\mathfrak{a}}
\newcommand{\hh}{\mathfrak{h}}
\renewcommand{\gg}{\mathfrak{g}}
\newcommand{\bb}{\mathfrak{b}}
\newcommand{\pp}{\mathfrak{p}}
\newcommand {\BB} {\mathfrak {B}}
\newcommand {\DD} {\mathscr {D}}
\newcommand {\OO} {\mathscr {O}}
\newcommand {\MM} {\mathscr {M}}

\renewcommand{\phi}{\varphi}
\newcommand{\eps}{\epsilon}
\newcommand{\del}{\partial}
\renewcommand{\mod}{\mathrm{mod~}}

\newcommand{\doublebrace}[4]{\left\{\begin{array}{ll}#1 & #2 \\ #3 & #4\end{array}\right.}
\newcommand{\triplebrace}[6]{\left\{\begin{array}{ll}#1 & #2 \\ #3 & #4\\ #5 & #6\end{array}\right.}
\newcommand{\half}{\frac{1}{2}}
\newcommand{\threehalfs}{\frac{3}{2}}
%\newcommand {\eqdef} {\stackrel {def}{=}}
\newcommand {\eqdef} {:=}
\newcommand {\eqdefrev} {=:}

\newcommand{\refle}[1]{Lemma \ref{#1}}
\newcommand{\refth}[1]{Theorem \ref{#1}}
\newcommand{\refcor}[1]{Corollary \ref{#1}}
\newcommand{\refse}[1]{Section~ \ref{#1}}
\newcommand{\refsec}[1]{Section~ \ref{#1}}
\newcommand{\refeq}[1]{(\ref{#1})}
%\theorembodyfont{\rm}
%-----------------document specific commands----------------%
\newcommand {\univ} {U}
\newcommand {\gkdim} {\mathrm{GKdim }}
\newcommand {\gkm} {$(\gg,\kk)$-module}
\newcommand {\kkm} {$(\kk,\kk)$-module}
\newcommand{\Fou}{\mathrm{Fou}}
%---------------------------------%


\begin{document}

\bibliographystyle{amsplain}

\ifx\useHugeSize\undefined
\else
\Huge
\fi

\relax

\title{ Bounded Simple $\left(\mathfrak g,\sl(2)\right)$-modules for $\rk \mathfrak g=2$}

\author{Ivan Penkov}
\address{
Jacobs University Bremen \\
School of Engineering and Science,
Campus Ring 1,
28759 Bremen, Germany}
\email{i.penkov@jacobs-university.de}

\author{Vera Serganova}
\address{Department of Mathematics, University of California Berkeley, Berkeley
  CA 94720, USA}
\email{serganov@math.berkeley.edu}

\subjclass[2000]{Primary 17B10, Secondary 22E46}
\begin{abstract}
This paper is a continuation of our work [PS2] in which we prove some general results about simple $\left(\mathfrak g, \mathfrak k \right)$-modules with bounded $\mathfrak k$-multiplicities (or bounded simple $\left(\mathfrak g, \mathfrak k \right)$-modules). In the absence of a classification of bounded simple $\left(\mathfrak g, \mathfrak k \right)$-modules in general, it is important to understand some special cases as best as possible. Here we consider the case $\mathfrak k=\sl(2)$. It turns out that in order for an infinite-dimensional bounded simple $\left(\mathfrak g, \sl(2) \right)$-module to exist, $\mathfrak g$ must have rank 2, and, up to conjugation, there are five possible embeddings $\sl(2)\rightarrow \mathfrak g$ which yield infinite-dimensional bounded simple $\left(\mathfrak g, \sl(2) \right)$-modules.

Our main result is a detailed description of the bounded simple $\left(\mathfrak g, \sl(2)\right)$-modules in all five cases. When $\gg \simeq \sl(2)\oplus\sl(2)$ we reproduce in modern terms some classical results from the 1940's. When $\gg \simeq \sl(3)$ and $\sl(2)$ is a principal subalgebra, bounded simple $\left( \sl(3), \sl(2)\right)$-modules are Harish-Chandra modules and our result singles out all Harish-Chandra modules with bounded $\sl(2)$-multiplicities. A case where the result is entirely new is the case of a principal $\sl(2)$-subalgebra of $\gg=\sp(4)$. 


\end{abstract}
\maketitle
\section{Introduction}

The classification of simple Harish-Chandra modules is a celebrated result and there is an extensive literature on the general topic of Harish-Chandra modules, see for instance [KV] and the references therein. Algebraically, Harish-Chandra modules are $\left(\gg, \kk \right)$-modules for a symmetric subalgebra $\kk$ of a semisimple Lie algebra $\gg$, and in the last decade an intense exploration of more general $(\gg, \kk)$-modules for not necessarily symmetric subalgebras $\kk$ has begun, [PS1],[PSZ], [PZ1], [PZ2], [PZ3]. A most notable result in this direction is the classification of simple $(\gg, \kk)$-modules of finite type and a generic minimal $\kk$-type carried out in [PZ2]. Nevertheless, the classification problem for $(\gg, \kk)$-modules of finite type with an arbitrary minimal $\kk$-type  is still open even when $\rk \gg=2$ and $\rk\kk=1$.

In the recent paper [PS2] we concentrated on the interesting subclass of bounded $(\gg, \kk)$-modules (the definition see in Section 1 below) and proved some general results regarding the existence of such modules.  In particular, we established sufficient and necessary conditions on a reductive in $\gg$ subalgebra $\kk$ for the existence of a simple infinite-dimensional bounded $(\gg, \kk)$-module. If $\kk \simeq \sl(2)$, simple $(\gg, \sl(2))$-modules of finite type exist for any simple $\gg$, see for instance [PSZ] or [PZ3]. It turns out however, that bounded $(\gg, \sl(2))$-modules are very special and exist only for $\rk\gg=2$. The classification of such modules is rather intriguing as they are the "smallest", and thus highly non-generic, $(\gg,\sl(2))$-modules.

This classification is carried out in the present paper. We show first that, up to conjugation, there are precisely five possibilities for embedding $\sl(2)$ into a Lie algebra of $\gg$ of rank 2 so that bounded infinite-dimensional $(\gg, \sl(2))$-modules exist:
 $\sl(2)$ as the diagonal subalgebra of $\sl(2)\oplus\sl(2)$, $\sl(2)$ as a root subalgebra or a principal $\sl(2)$ subalgebra of $\sl(3)$, and $\sl(2)$ as a root subalgebra corresponding to a short root or as a principal subalgebra of $\sp(4)$. 

We then give a classification and a detailed description (we compute characters and minimal $\sl(2)$-types) of all bounded $(\gg, \sl(2))$-modules. In the case when $\gg\simeq \sl(2)\oplus\sl(2)$ our results are just a modern reproduction of classical results, in all other cases they are new. The most interesting new case is that of a principal $\sl(2)$-subalgebra of $\sp(4)$.


\textbf{Acknowledgement.} This paper has been written in close contact
with Gregg Zuckerman who has supported us on several occasions with
valuable advice. David Vogan, Jr. has also generously shared his
knowledge of Harish-Chandra modules with us. We thank T. Milev for
reading the manuscript carefully and checking some of the
calculations. Finally, we acknowledge the hospitality and support of the Max Planck Institute for Mathematics in Bonn.

\section{General definitions and preliminary results}

The ground field is $\CC$. 

Let $\gg$ be a semisimple (finite-dimensional) Lie algebra and
$\kk\subset\gg$ be a reductive in $\gg$ subalgebra. $U(\cdotp)$ stands
for enveloping algebra, $U=U(\gg)$ and $Z_U$ is the center of $U$. A $(\gg, \kk)$\textit{-module}  $M$ is a $\gg$-module  $M$ on which $\kk$ acts locally finitely, i.e. $\dim U(\kk) \cdotp m <\infty$, $\forall m\in M$. A $(\gg, \kk)$-module $M$ has \textit{finite type over} $\kk$ if the Jordan--H\"{o}lder multiplicity of any fixed simple finite-dimensional $\kk$-module $V$ (such a $V$ is called a $\kk$\textit{-type}) in arbitrary finite-dimensional $\kk$-submodules of $M$ is bounded. A $(\gg, \kk)$-module is \textit{bounded} if the above multiplicities are bounded by a constant not depending on the $\kk$-type $V$. A reductive in $\gg$ subalgebra $\kk\subset \gg$ is \textit{bounded} if there exists an infinite-dimensional simple bounded $(\gg, \kk)$-module $M$.
A bounded subalgebra $\kk\subset\gg$ is \textit{strictly bounded} if there is a simple bounded $(\gg,\kk)$-module $M$ on which no simple ideal of $\gg$ acts locally finitely. The following necessary conditions on a subalgebra $\kk$ to be bounded, or strictly bounded, are proved in [PS2] (Theorem 4.1 and Corollary 4.6).

\begin{theorem}\label{1}
 Let $\kk$ be a bounded reductive subalgebra of a semisimple Lie algebra $\displaystyle \gg=\oplus_i \gg_i$ ($\gg_i$ being the simple ideals of $\gg$).


a) If $M$ is a simple bounded $(\gg, \kk)$-module and the algebra of $\kk$-invariants $\gg_{i_0}^{\kk}$ is not abelian for some $i_0$, then $M \simeq M_{i_0}\otimes \overline{M_{i_0}}$, where $M_{i_0}$ is a simple finite-dimensional $\gg_{i_0}^{\kk}$-module and $\overline{M_{i_0}}$ is a simple bounded $(\oplus_{i\neq {i_0}}\gg_i,(\oplus_{i\neq {i_0}}\gg_i)\bigcap\kk)$-module.


b) If $r_\gg$ is the half-dimension of a nilpotent orbit of minimal positive dimension in $\gg$, then
\begin{equation}\label{1}
r_{\gg}\leq b_{\kk},
\end{equation}
where $b_\kk$ is the dimension of a Borel subalgebra of $\gg$.



c) If $\kk$ is strictly bounded, then

$$\sum_i r_{\gg_i}\leq b_{\kk}.$$



\end{theorem}

In [PS2] we also established the following sufficient condition for a reductive in $\gg$ subalgebra $\kk\subset\gg$ to be bounded. Recall that a finite-dimensional module $W$ over an algebraic group $H$ is \textit{spherical} if a Borel subgroup $B_{H}$ has an open orbit in $W$. 

\begin{theorem}\label{2}
Let $K\subset G\subset GL(V)$ be a chain of reductive algebraic
groups, and let $V'\subset V$ be a 1-dimensional space whose
stabilizers in $G$ and $K$ are parabolic subgroups $P\subset G$ and
$Q\subset K$. Then, if $(V')^*\otimes(\gg\cdot V'/\kk\cdot V')$ is a
spherical module over a reductive part $Q_0$ of $Q$, $\kk$ is a bounded subalgebra of $\gg$.
\end{theorem} 


\section{Bounded subalgebras of a rank-2 Lie Algebra}\label{secRank2Case}

Our main interest in this paper are infinite-dimensional bounded $(\gg, \kk)$-modules for $\kk\simeq\sl(2)$. \refth{1} implies that if $\kk\simeq\sl(2)$ is a strictly bounded subalgebra of $\gg=\oplus_i \gg_{i}$, then $\sum_i r_{\gg_i}\leq 2$. This is easily seen to imply $\rk \gg=2$. Therefore, in the rest of the paper we restrict ourselves to the case when $\rk\gg=2$. The following theorem classifies more generally all reductive in $\gg$ bounded subalgebras $\kk\subset\gg$ under the assumption that $\rk\gg=2$.


\begin{theorem}\label{thPairsList}
Let $\gg$ be a semisimple Lie algebra of rank 2 and $\kk\subset\gg$ be a reductive in $\gg$ bounded subalgebra. The following is a complete list of such pairs up to conjugation by inner automorphisms.
\begin{itemize}
\item[(1)] $\gg\simeq\sl(2)\oplus\sl(2)$: $\kk\simeq \gl(2)$ is a direct sum of a simple ideal and a Cartan subalgebra of the other simple ideal,
$\kk\simeq\sl(2)$ is a
  diagonal subalgebra, or $\kk$ is any non-trivial toral subalgebra;
\item[(2)] $\gg\simeq\sl(3)$: $\kk$ is a root subalgebra isomorphic to $\sl(2)$ or $\gl(2)$, $\kk$ is a principal   $\sl(2)$-subalgebra, or $\kk$ is a Cartan subalgebra;
\item[(3)] $\gg\simeq\sp(4)$: $\kk \simeq \sl(2)\oplus \sl(2)$ is the subalgebra generated by the long roots, $\kk \simeq \gl(2)$ is any root subalgebra, $\kk\simeq \sl(2)$ is a root   subalgebra corresponding to a short root, $\kk$ is a principal $\sl(2)$-subalgebra, or $\kk$ is a Cartan subalgebra;
\item[(4)] $\gg\simeq G_2$: $\kk$ is any subalgebra containing a Cartan subalgebra, in this case $\kk\simeq \sl(3)$, $\kk\simeq\sl(2)\oplus\sl(2)$, or $\kk\simeq \gl(2)$.
\end{itemize}
\end{theorem}
\begin{proof} The inequality \refeq{1} implies that a 1-dimensional toral subalgebra is not bounded in all cases but (1). In (1) any 1-dimensional
  toral subalgebra $\mathfrak {t}$ is bounded as the outer tensor
  product of a Verma module over a suitable ideal of $\gg$ with the
  trivial module of the complementary ideal of $\gg$ is always bounded as a
  $(\gg,\mathfrak t)$-module.

Similarly, \refeq{1} implies that a Cartan subalgebra is not bounded in $G_2$. In all other cases it is well known to be bounded, see for instance \cite{F}.

If $\kk \simeq \sl(2)$ then $\kk$ is not bounded in $G_2$ again by
\refeq{1}, and if $\kk$ is an ideal of $\gg=\sl(2)\oplus \sl(2)$, it
is not bounded by  \refth{1} a). Furthermore, if $\kk \simeq \sl(2)$ is
a root subalgebra of $\gg=\sp(4)$ corresponding to a long root, then
$\kk$ is not bounded by \refth{1} a). For the remaining five possible embeddings of $\sl(2)$ into a Lie algebra of rank 2, the image
$\kk$ is always a bounded subalgebra. This follows for instance from
the explicit description of bounded $(\gg,\kk)$-modules which we present in
Sections 4-7 of this paper.

For any embedding of $\gl(2)$ into a Lie algebra $\gg$ of rank 2,
$\gg\ncong G_2$, any generalized Verma module, corresponding to
a parabolic subalgebra $\pp$ which contains the image $\kk$ of
$\gl(2)$, is a bounded $(\gg,\kk)$-module.

Consider next the case $\kk\simeq \sl(2)\oplus \sl(2)\subset \gg$ for $\gg=\sp(4)$
or $G_2$. Here the pair $(\gg,\kk)$ is symmetric. In
\cite{V1} and \cite{V3} ladder $(\gg,\kk)$-modules are
constructed. Fix a Borel subalgebra $\bb_{\kk}\subset \kk$. By
definition, a ladder module $M$ has the $\kk$
decomposition $M=\bigoplus_{n\in \ZZ_{\geq 0}}V_{\mu + n\beta}$, where
$\mu$ is some integral $\bb_{\kk}$-dominant weight and $\beta$ is
the $\bb_{\kk}$-highest weight of $\gg/{\kk}$. Clearly, a ladder
module is multiplicity-free and hence bounded. Moreover, it remains bounded with respect to any
$\gl(2)$-subalgebra of $\kk$. Hence any image of $\gl(2)$ in $\sp(4)$
or $G_2$ is bounded.

The only remaining case is $\gg=G_2, \kk\simeq \sl(3)$. To show that $\kk$ is bounded we use \refth{2} with $V$ being the 7-dimensional $G_2$-module. Then as a $\kk$-module $V$ is isomorphic to $V_{\omega_1}\oplus V^*_{\omega_1}\oplus \CC$. One can fix a Borel subalgebra $\bb\subset \gg$ so that there exists a $\bb$-stable one-dimensional subspace $V'\subset
V_{\omega_1}^*$. Then $Q_0\simeq \GL(2)$ and
$$(V')^*\otimes (\gg\cdot V'/\kk\cdot V')\simeq \Lambda^2(E)\otimes (E^*\oplus \mathbb C)$$
where $E$ is the standard $GL(2)$-module. It is easy to check that it 
is a spherical $Q_0$-module.
\end{proof}

In the rest of this paper $\gg$ will be of rank 2, and $\kk$ will be isomorphic to $\sl(2)$. By $V_k$ we denote the $k+1-$dimensional $\kk$-module, and we write $c(M)$ for the $\kk$-character of any  $(\gg,\kk)$-module $M$ of finite type over $\kk$:
\[
c(M):=\sum_{k\geq 0} (\dim M^k)z^k,
\]
where $M^k=\Hom_\kk(V_k,M)$.
By definition, $c(M)$ is a formal power series in $z$. The
\emph{minimal $\kk$-type} of $M$ is $V_t$ where $t\in\ZZ_{\geq 0}$ is
minimal with $M^t\neq 0$. A $(\gg,\kk)$-module of finite type $M$ is
\emph{even} (respectively, \emph{odd}) if $M^t=0$ for all $t\in 1+2\ZZ$ (resp., $t\in 2\ZZ$).

Let $\CC((z))$ be the algebra of Laurent series and $\CC ((z))'$ be the span of vectors in $\CC((z))$ of the form $z^j+z^{-j-2}$ for $j\in\ZZ$ ($\CC((z))'$ is not a subalgebra). Note that $\CC((z))'$ is a complement to the subspace $\CC[[z]]$ of $\CC((z))$. In what follows we denote by $\pi$ the projection onto the second summand in the direct sum $\CC((z))=\CC((z))'\oplus\CC[[z]]$, and we set $z^p\otimes z^q:=\sum_{0\leq k\leq q}z^{p+q-2k}$ for $p\geq q$ and $z^p\otimes z^q:=z^q\otimes z^p$ for $p<q$.

\begin {lemma}\label{le202}~
\begin{itemize}
\item[(a)] For any $f(z)\in \CC((z))$ and any $j\in \ZZ$,  $\pi(f(z)(z^j+z^{-j}))=\pi(\pi(f(z)(z^j+z^{-j})))$.
\item[(b)] For any $(\kk,\kk)$-module $M$ of finite type over $\kk$
\[
c(M\otimes V_i)=\pi(c(M)\sum_{0\leq k\leq i}z^{i-2k}),
\]
for all $i\in \NN$.
\end{itemize}
\end{lemma}
\begin {proof}~
\begin {itemize}
\item [(a)] It suffices to check that for any $\psi(z)\in C((z))'$, $\psi(z)(z^j+z^{-j})\in \CC((z))'$, and this is obvious.
\item [(b)] It suffices to check that, for any $s\in \ZZ_{\geq 0}$
\[
\pi(z^s\otimes(\sum_{0\leq k\leq i}z^{i-2k}))=\sum_{0\leq k\leq \frac{|i-s|}{2}}z^{s+i-2k},
\]
which is also obvious.
\end{itemize}
\end{proof}

Finally, by $\Gamma_\kk$ we denote the functor of $\kk$-finite vectors:
\[
\Gamma_\kk:\gg-\text{mod}\leadsto(\gg,\kk)-\text{mod},
\]
\[
M\mapsto\{m\in M|\dim(U(\kk)\cdot m)<\infty\}.
\]

\section{Classification and $\kk$-characters of simple $(\sl(2)\oplus\sl(2),\sl(2))$-modules}\label{secKchars}

\refth{thPairsList} singles out the cases when $\kk\simeq\sl(2)$ is a bounded subalgebra of a rank-2 Lie algebra. The simplest case is when $\gg=\sl(2)\oplus\sl(2)$ and $\kk\subset\gg$ is the diagonal subalgebra. Here all simple $(\gg,\kk)$-modules are bounded and are moreover multiplicity-free. This follows, for instance, from
the algebraic subquotient theorem, see \cite{Dix}, Ch.~9. These $(\gg,\kk)$-modules are historically among the first examples of $(\gg,\kk)$-modules studied. They have been classified already in 1947 by Gelfand and Naimark
\cite{GN} and by Bargmann \cite{B}, and have been constructed also by Harish-Chandra
around the same time, \cite{HC}.
A fundamental more modern and much more general
reference is the article \cite{BG}, where however this explicit
example is not written in detail.
In the present section we give a quick self-contained description of
all simple $(\gg,\kk)$-modules based on the approach of \cite{BG}.

\begin{lemma}\label{lmc}
Let $\Omega_1, \Omega_2\in U(\gg)$ be the Casimir elements of the two $\sl(2)$-direct summands of $\gg$, and $\Omega\subset U(\kk)\subset U(\gg)=U$ be the Casimir element of $\kk$. Then $\Omega_1, \Omega_2$ and $\Omega$ generate $U(\gg)^{\kk}$.
\end{lemma}
\begin{proof}
Straightforward computation. A more general result is proved by F. Knop in \cite{Kn1}.
\end{proof}

\begin{corollary}
Every simple $(\gg,\kk)$-module is multiplicity-free.
\end{corollary}

Denote by $\chi(a,b)$ the central character of the Verma $\gg$-module
with highest weight $(a_1,b_1)$, where the notation $(c,d)$ is shorthand for the weight $c\omega_{\mathrm{left}}+d\omega_{\mathrm{right}}$, $\omega_{\mathrm{left}}$ (respectively, $\omega_{\mathrm{right}}$) being the fundamental weight of the first (respectively, second) direct summand of $\gg$.



\begin{lemma}\label{lmch}
If $V_n$ is the minimal $\kk$-type of a simple infinite-dimensional $(\gg,\kk)$-module $M$, then
\begin{equation}\label{eqChar1}
c(M)=z^n+z^{n+2}+z^{n+4}+\dots\quad .
\end{equation}
\end{lemma}
\begin{proof}
To prove \refeq{eqChar1} it suffices to show that $V_n$, $V_{n+2}$,
$V_{n+4}$, etc. are precisely all $\kk$-types of $M$. The absence of
other $\kk$-types follows from the fact that as a $\kk$-module $\gg$
is isomorphic to $V_2\oplus V_2$, hence when acting by $\gg$ on
$V_{n+2i}$ one can only obtain $\kk$-constituents of $(V_2\oplus
V_2)\otimes V_{n+2i}$, i.e. $V_{n+2(i-1)}$, $V_{n+2i}$ and
$V_{n+2(i+1)}$. To show that for each $i>0$ $V_{n+2i}$ is a
$\kk$-constituent of $M$, note that if $V_{n+2i}$ were not a
constituent of $M$, then  when acting by $\gg$ on $V_{n+2(i-t)}$ for $t\geq 1$ one would not be able to obtain a constituent of the from $V_{n+2(i+r)}$ for $r\geq 1$. Hence $M$ would turn being finite-dimensional, a contradiction.
\end{proof}

\begin{lemma}
Let $M$ be a simple $(\gg,\kk)$-module with minimal $\kk$-type $V_0$. Then the central character of $M$ equals $\chi(a,a)$ for some $a\in \CC$.
\end{lemma}
\begin{proof}
Since $\gg\simeq \kk\oplus\kk$, the $\gg$-module  $U\otimes_{U(\kk)}V_0$ is isomorphic to $U(\kk)$. The latter is endowed with a $U\simeq U(\kk)\otimes U(\kk)$-module structure via left multiplication by elements of $U(\kk)\otimes 1$ and right
 multiplication by elements of $1\otimes U(\kk)$.
Moreover, the action of $\Omega_1$ and $\Omega_2$ coincides on $U(\kk)$. Since $M$ is a quotient of the $\gg$-module $U(\kk)$, the action of $\Omega_1$ and $\Omega_2$ coincides on $M$, hence the Lemma.
\end{proof}
\begin{lemma}
Let $M$ be a simple $(\gg,\kk)$-module. Then the central character of $M$ equals $\chi (a,a+n)$ for some $a\in \CC$ and some $n\in\ZZ$. Moreover, the parity of $n$ equals the parity of $k$ where $V_k$ is the minimal $\kk$-type of $M$.
\end{lemma}
\begin{proof}
Let $M$ have central character $\chi(\alpha,\beta)$.
Consider the $\gg$-module $M\otimes (V_0\boxtimes V_k)$, where
the $\gg=\kk\oplus\kk$-module
$V_0\boxtimes V_k$ is endowed with a $\gg$-module structure via
the isomorphism $\gg\simeq \kk\oplus\kk$. Then
$\hom_\kk(V_0,M\otimes(V_0\boxtimes V_k))\neq 0$, hence a simple
subquotient of $M\otimes (V_0\boxtimes V_k)$ has central character
$\chi (a,a)$ for some $a$. On the other hand, the central characters
of all simple subquotients of $M\otimes(V_0\boxtimes V_k)$ are of the
form $\chi(\alpha,\beta-n)$ for $n$ running over the set of weights of
$V_k$.
Therefore $\alpha=a$, $\beta-n=a$, i.e. the Lemma follows.
\end{proof}

\begin{lemma}\label{le00}
For any central character $\chi$, up to isomorphism there is at most one infinite
dimensional simple $(\gg,\kk)$-module with this central character.
\end{lemma}
\begin{proof}
Let $M',M''$ be two simple $(\gg,\kk)$-modules with central character
$\chi$. Then, by \refle{lmch}, for some $m$
$\hom_\kk(V_m,M')=\hom_\kk(V_m,M'')=\CC$. Therefore $M'$ and $M''$ are
isomorphic to simple quotients of the $\gg$-module $U\otimes_{Z_U
  U(\kk)}V_m$, where $Z_U$ acts on $V_m$ via the central character
$\chi$. The fact that $U^\kk\subset Z_U U(\kk)$ (\refle{lmc}) implies
that $\hom_\kk(V_m, U\otimes_{Z_U U(\kk)}V_m)=\CC$ for every $m\geq
0$. Hence $U\otimes_{Z_U U(\kk)}V_m$ has a unique proper maximal
submodule, and in this way also a unique simple quotient. Therefore $M'\simeq M''$.
\end{proof}

In the rest of this section  we only consider central characters of
the form $\chi(a,a-n)$ for $n\in\mathbb Z_{\geq 0}$. If $a\in \ZZ$, we assume in addition that $a\geq 0$ and $a-n \leq 0$. By $M_c$ denote the Verma module over $\kk$ with highest weight $c-1$. Note that
for $a,a-n$ as above, $\hom_{\CC}(M_{a},M_{a-n})$ is a $\gg$-module with central character $\chi(a,a-n)$. Define
\[
W_{a,a-n}:=\Gamma_{\kk}(\hom_{\CC} (M_{a},M_{a-n})).
\]
\begin{theorem}\label{thMainSCT}~
\begin{itemize}
\item[(a)] Fix $a\in\CC\backslash \ZZ_{< 0}$ and $n\in \ZZ_{\geq 0}$ such that $a-n\leq 0$ for integer $a$. The $\gg$-module $W_{a,a-n}$ is the unique (up to isomorphism) simple infinite-dimensional $(\gg,\kk)$-module with central character $\chi(a,a-n)$.
\item[(b)] $c(W_{a,a-n})=z^n+z^{n+2}+z^{n+4}+\dots\quad .$
\end{itemize}
\end{theorem}
\begin{proof} Note that to compute the
$\kk$-character of
$\Gamma_{\kk}(\hom_{\CC}(M_a,M_{a-n}))$ it suffices to compute
  $\hom_\kk(V_m,\hom_{\CC}(M_a,M_{a-n}))$ for all $m\in \ZZ_{\geq 0}$. However,
\[
\hom_\kk(V_m,\hom_{\CC}(M_a,M_{a-n}))=\hom_\kk(M_a,M_{a-n}\otimes V_m^*),
\]
and
\[
\hom_\kk(M_a, M_{a-n}\otimes V_m^*)=\left\{\begin{array}{cc}\CC&\text{for~}m-n\in2\ZZ_{\geq 0}\\0&\text{otherwise}\end{array}\right. .
\]
Hence
\[
c(W_{a,a-n})=z^n+z^{n+2}+z^{n+4}+\dots\quad .
\]
The simplicity of $W_{a,a-n}$ follows from the observation that if simple, $W_{a,a-n}$ would have a finite-dimensional subquotient, but
there is no finite-dimensional $\gg$-module with central character
$\chi(a,a-n)$ for $a\in \CC\backslash \ZZ$ or $a=0$.  If $a\in
\ZZ $, the finite-dimensional $\gg$-module with central
character $\chi(a,a-n)$ is isomorphic to $V_{a-1}\boxtimes V_{n-a-1}$
whose $\kk$-character is $z^{n-2}+z^{n-4}+...+z^{|n-2a-2|}$, and
hence it can not be a subquotient of $W_{a,a-n}$.
\end{proof}

\section{Classification and $\kk$-characters of simple bounded $(\sl(3),\sl(2))$-modules}\label{secKchars3-2}

Throughout this section $\gg=\sl(3)$ and $\kk\simeq\sl(2)\subset\gg$.

\subsection{The root case.}\label{subsecTheRootCase}
In this subsection we fix a Cartan subalgebra $\hh\subset\gg$ and simple roots $\alpha_1,\alpha_2\in\hh^*$ which define a Borel subalgebra $\bb^+\subset\gg$. We also fix $\kk$ to be the $\sl(2)$-subalgebra generated by the root spaces $\gg^{\pm\alpha_1}$. There are two parabolic subalgebras containing $\kk$ and $\hh$: $\pp^+:=(\hh+\kk)\oplus\gg^{\alpha_2}\oplus\gg^{\alpha_1+\alpha_2}$, $\pp^-:=(\hh+\kk)\oplus\gg^{-\alpha_2}\oplus\gg^{-\alpha_1-\alpha_2}$. Note that $\bb^+\subset\pp^+$ and define $\bb^-$ to be the Borel subalgebra with simple roots $\alpha_1,-\alpha_1-\alpha_2$. Then $\bb^-\subset\pp^-$. In addition, we fix generators $h_i\in[\gg^{\alpha_i},\gg^{-\alpha_i}]$ and denote by $\omega_i$, for $i=1,2$, the corresponding dual basis of $\hh^*$. Then $\rho_{\bb^+}=\omega_1+\omega_2$, $\rho_{\bb^-}=\omega_1-2\omega_2$.

\begin{lemma}\label{le81}
Let $M$ be a simple bounded infinite-dimensional $(\gg,\kk)$-module. Then $\gg[M]=\pp^\pm$.
\end{lemma}
\begin{proof}
Since $\hh\subset\gg^\kk\oplus\kk$, we have $\hh\subset\gg[M]$. Put $M_0:=\{m\in M| \gg^{\alpha_1}\cdot m=0\}$ and choose generators $x$ and $y$ of the respective root spaces $\gg^{-\alpha_2}$ and $\gg^{\alpha_1+\alpha_2}$. A straightforward computation shows that for any $i,j\in \ZZ_{\geq 0}$, $(x^iy^j)\cdot v\in M_0$ if $v$ is any non-zero vector in $M_0$ such that $h_1\cdot v= \nu(h_1) v$ for some $\nu\in(\hh\cap\kk)^*$. Therefore the assumption that $x,y\notin \gg[M]$ implies that the multiplicity of $V_{\nu+i+j}$ is at least $i+j$, which contradicts the boundedness of $M$. Hence $\gg^{-\alpha_2}\in\gg[M]$ or $\gg^{\alpha_1+\alpha_2}\in\gg[M]$, and consequently $\gg[M]=\pp^\pm$.
\end{proof}

Let $F_{a,b}^\pm$ be the simple finite-dimensional $\pp^\pm$-module with $\bb^\pm$-highest weight $a\omega_1+b\omega_2$. Define $L^\pm_{a,b}$ as the unique simple quotient of $U(\gg)\otimes_{U(\pp^\pm)}F^\pm_{a,b}$. Then $L^\pm_{a,b}$ are bounded $(\gg,\kk)$-modules, and the existence of an isomorphism $L^\pm_{a,b}\simeq L^\mp_{a',b'}$ implies $\dim L_{a,b}^\pm<\infty$.

\begin{theorem}\label{th82}
Let, as above, $\kk\simeq \sl(2)$ be a root subalgebra of $\gg=\sl(3)$.
\begin{itemize}
\item[(a)] Any infinite-dimensional bounded $(\gg,\kk)$-module is isomorphic either to $L_{a,b}^+$ for $a\in\ZZ_{\geq 0}$, $b\in\CC\backslash\ZZ_{\geq 0}$ or to $L^-_{a,b}$ for $a\in\ZZ_{\geq 0}$, $-a-b\in\CC\backslash\ZZ_{\geq 0}$.
\item[(b)]
\begin{equation}\label{eq81}
c(L_{a,b}^\pm)=1+2z+\dots+az^{a-1}+(a+1)(z^a+z^{a+1}+\dots)
\end{equation}
for all $a\geq 0$ and for those $b$ which do not satisfy the conditions $-b\in\ZZ_{\geq 2}$, $a+b\in\ZZ_{\geq -1}$ for $L_{a,b}^+$, and respectively the conditions $a+b\in\ZZ_{\geq 2}$, $-b\in\ZZ_{\geq-1}$ for $L_{a,b}^-$.
\item[(c)] If $-b \in \ZZ_{\geq 2}$, $a+b\in\ZZ_{\geq -1}$, then
\begin{equation}\label{eq82}
c(L_{a,b}^+)=z^{-b-1}+2z^{-b}+\dots+(a+b+1)z^{a-1}+(a+b+2)(z^a+z^{a+1}+\dots),
\end{equation}
and if $a+b\in\ZZ_{\geq 2}$, $-b\in \ZZ_{\geq -1}$, then
\begin{equation}\label{eq83}
c(L_{a,b}^-)=z^{a+b-1}+2z^{a+b}+\dots+(1-b)z^{a-1}+(2-b)(z^a+z^{a+1}+\dots).
\end{equation}
\end{itemize}
\end{theorem}
\begin{proof}
Let $M$ be a simple infinite-dimensional bounded
$(\gg,\kk)$-module. Then, by \refle{le81}, $\gg[M] =\pp^\pm$. If
$\gg[M]=\pp^+$, let $M^+$ be a simple finite-dimensional
$\pp^+$-submodule of $M$. Then $M^+\simeq F^{+}_{a,b}$ for some
$a\in\ZZ_{\geq 0}$ and some $b\in \CC$, and there is an obvious
surjection of $\gg$-modules $U(\gg)\otimes_{U(\pp^+)}F^+_{a,b}\to
M$. Hence $M$ is isomorphic to the unique simple quotient $L_{a,b}^+$
of $U(\gg)\otimes_{U(\pp^+)} F_{a,b}^+$. However, $L_{a,b}^+$ is
finite-dimensional iff $b\in\ZZ_{\geq 0}$, therefore (a) follows for
the case when $\gg[M]=\pp^+$. The case $\gg[M]=\pp^-$ is obtained by replacing $b$ with $-a-b$ which corresponds to the replacement of the simple root $\alpha_2$ of $\bb^+$ by the simple root $-\alpha_1-\alpha_2$ of $\bb^-$.

Statements (b) and (c) follow from a non-difficult reducibility
analysis for the induced module $U(\gg)\otimes_{U(\pp^\pm)} F_{a,b}^\pm$. Note first of all that $\ch_\kk(U(\gg)\otimes_{U(\pp^\pm)} F_{a,b}^\pm)$ is always given by the right-hand side of \refeq{eq81}. Indeed as $\kk$-modules $\gg/\pp^\pm$ and $F_{a,b}^\pm$ are isomorphic respectively to $V_1$ and $V_a$, therefore
\[
c(U(\gg)\otimes_{U(\pp^\pm)} F_{a,b}^\pm)=c(S^\cdot(V_1)\otimes V_a).
\]
A straightforward computation shows that $c(S^\cdot(V_1)\otimes V_a)$ is nothing but the right hand side of \refeq{eq81}.

We claim now that $U(\gg)\otimes_{U(\pp^\pm)} F_{a,b}^\pm$ is irreducible precisely when $b$ does not satisfy the respective conditions stated in (b). Consider first the case of $\pp^+$. Then $U(\gg)\otimes_{U(\pp^+)} F_{a,b}^+$ is irreducible if and only if there exists $w\in W\backslash W_\kk$ such that
\begin{equation}\label{eq84}
(w((a+1)\omega_1+(b+1)\omega_2)-(\omega_1+\omega_2))(h_1)\in \ZZ_{\geq 0}
\end{equation}
and
\begin{equation}\label{eq85}
(w((a+1)\omega_1+(b+1)\omega_2)-(\omega_1+\omega_2))=a\omega_1+b\omega_2-m_1\alpha_1-m_2\alpha_2
\end{equation}
for some $m_1,m_2\in \ZZ_{\geq 0}$. The only non $\bb^+$-dominant solution of \refeq{eq84} and \refeq{eq85} is $w=w_{\alpha_1+\alpha_2}$ and $-b\in \ZZ_{\geq 2}, a+b\in\ZZ_{\geq -1}$. Moreover, in the latter case $L_{a,b}^+\simeq(U(\gg)\otimes_{U(\pp^+)} F_{a,b}^+)/L_{-b-2,-a-2}^+$, where $c(L^+_{-b-2,-a-2})$ is given by the right hand side of \refeq{eq81} with $a$ replaced by $-b-2$. An immediate computation shows that $c( L_{a,b}^+)$ is given in this case by the right hand side of \refeq{eq82}, therefore (b) and (c) are proved for the case of $\pp^+$. The case of $\pp^-$ is obtained by interchanging the parameter $b$ in \refeq{eq82} with $-a-b$.
\end{proof}

\begin{corollary}
Let $\gg$ and $\kk$ be as above.
\begin{itemize}
\item[(a)] The minimal $\kk$-type of a simple bounded infinite-dimensional $(\gg,\kk)$-module can be arbitrary. The multiplicity of the minimal $\kk$-type is always 1.
\item[(b)] The following is a complete list of multiplicity-free simple infinite-dimensional $(\gg,\kk)$-modules:
\begin{itemize}
\item $L_{0,b}^+$ for $b\in\CC\backslash \ZZ_{\geq 0}$,
\item $L_{0,b}^-$ for $-b\in\CC\backslash \ZZ_{\geq 0},$
\item $L_{a,b}^+$ for $a+b=-1$, $-b\in \ZZ_{\geq 2}$,
\item $L_{a,b}^-$ for $b=1$, $a+b\in \ZZ_{\geq 2}$.
\end{itemize}
\end{itemize}
\end{corollary}

\subsection{The principal case.}
Let now $\kk$ be a principal $\sl(2)$-subalgebra of $\gg=\sl(3)$. 
The pair $(\gg,\kk)$ is well known to be symmetric and 
the simple $(\gg,\kk)$-modules have been studied extensively, see for
instance \cite{Fo} and \cite{Sp}. In principle one should be able to
identify all simple bounded modules in the known classification of
simple Harish-Chandra modules. However, we propose an alternative
approach which leads directly to all bounded simple
$(\gg,\kk)$-modules and their $\kk$-characters.  This is the first
case in which the richness of the theory of bounded (generalized)
Harish-Chandra modules becomes apparent.


We keep the notations $\hh,\bb^+, \alpha_1, \alpha_2$ from Subsection
\ref{subsecTheRootCase}. By $L_{a,b}$ we denote the simple
$\gg$-module with $\bb^+$-highest weight
$(a-1)\omega_1+(b-1)\omega_2$, by $V_{p,q}$ we denote the simple
finite-dimensional $\gg=\sl(3)$-module with $\bb^+$-highest weight
$p\omega_1+q\omega_2$ ($p,q\in\ZZ_{\geq 0}$), and $\chi(a,b)$ stands
for the central character of $L_{a,b}$. By $A$ we denote the Weyl algebra in the indeterminates $t,x,y$.

We first describe the primitive ideals of all simple bounded
$(\gg,\kk)$-modules. Let $\mathrm{GKdim}M$ denote the Gelfand-Kirillov
dimension of a $\gg$-module $M$ and $X_M$ denote the associated
variety of $M$.

\begin{lemma}\label{le100sl}
Let $M$ be an infinite-dimensional bounded simple $(\gg,\kk)$-module.
Then $\Ann M=\Ann L_{a,b}$, where $\dim L_{a,b}=\infty$, $a\in\ZZ_{>0}$, $b\in \ZZ_{>0}$ or $a+b\in\ZZ_{>0}$.
\end{lemma}
\begin{proof}
By Duflo's Theorem $\Ann M=\Ann L_{a,b}$ for some $a,b$. By Theorem
4.4 in [PS2],  
$\mathrm{GKdim}M\leq 2$. Since
$\mathrm{GKdim}M\geq\frac{1}{2}\mathrm{dim}X_M$ and
$\mathrm{GKdim}L_{a,b}=\frac{1}{2}\mathrm{dim}X_M$, we have
$\mathrm{GKdim}L_{a,b}\leq 2$. A straightforward computation shows that this latter
condition is equivalent to the condition on $(a,b)$ in the statement
of the Lemma.
\end{proof}

Let $\BB_\kk^{\chi}$ be the category of bounded $(\gg,\kk)$-modules which afford the central character $\chi$, see [PS2], Section 4.

\begin{corollary}\label{corPrimIdeal}
If $\BB^\chi_\kk$ is not empty, then $\chi=\chi(u+1-n,n+1)$ for some $n\in\ZZ_{\geq 0}$, where $u\in\CC\backslash \ZZ_{<n-1}$ or $u=-2$.
\end{corollary}

Note that the natural embedding of $\gl(3)$ into $A$ maps the center of
$\gl(3)$ to the line $\CC \EE$ for
$\EE:=t\partial_t+x\partial_x+y\partial_y$, and that the adjoint
action of the central element $\EE$ on $A$ defines a $\ZZ$-grading
$A:=\bigoplus _{i\in \ZZ}A_i$. Let $u\in\mathbb C$. Define the (associative) algebra $D^u$ as the
quotient of $A_0$ by the ideal generated by $\EE-u$. The embedding of
$\gg\to A_0$ induces a surjective homomorphism $\gamma_u: U(\gg)\to
D^u$. It is not difficult to show that if $u\in\ZZ$, $D^u$ is
isomorphic to the algebra of globally defined differential
endomorphisms of the line bundle $\OO_{\PP^2}(u)$ ($\PP^2$ being the
projective space with homogeneous coordinates $(x,y,z)$).

\begin{lemma}\label{le1021sl} Consider $D^u$ with its adjoint $\gg$-module structure. Then
\[
D^u\simeq \bigoplus_{m\geq 0}V_{m\rho}.
\]
\end{lemma}
\begin{proof} Let $\CC =A^0\subset A^1 \subset \dots \subset A$
  denote the standard filtration of $A$. A direct computation shows
  that as a $\gg$-module $A_0^m/A_0^{m-1}$ is
  isomorphic to
\[
V_{m,0}\otimes V_{0,m}=\oplus_{k=0}^m V_{k \rho}. 
\]
 After factorization by $\EE-u$, one obtains
\[
(D^u)^m/(D^u)^{m-1}\simeq V_{m \rho}. 
\]
\end{proof}

It is not difficult to see that the restriction of $\gamma_u$ to $U(\kk)$ is injective. Slightly abusing notation we identify $U(\kk)$ with its image in $D^u$. We will use the following expression for the standard basis $E,H,F$ of $\kk$:
\begin{equation}\label{princ}
E=t\partial_x+x\partial_y, H=2t\partial_t-2y\partial_y, F=2x\partial_t+2y\partial_x.
\end{equation}

\begin{lemma}\label {le102sl} The centralizer of $\kk$ in $D^u$ coincides with the center of $U(\kk)\subset D^u$.
\end{lemma}
\begin{proof}
As $V_{m\rho}^{\kk}=0$ for odd $m$ and  $V_{m\rho}^{\kk}=\CC$ for even
$m$ it is clear that the centralizer of $\kk$ in $D^u$ is generated by
the quadratic Casimir element $\Omega \in V_{2\rho}^{\kk}$.
\end{proof}

\begin{corollary}\label{cor98}
Every simple $(D^u,\kk)$-module is multiplicity-free. For any non-negative
 $m$, there exists at most one (up to isomorphism) 
simple $(D^u,\kk)$-module $M$ with $\hom_{\kk}(V_m,M)\neq 0$.
\end{corollary}
\begin{proof} It is well known that if $M$ is a simple $(\gg, \kk)$-module, then $M^{V}=\Hom_\kk(V, M)$ is a simple $U(\gg)^{\kk}$-module for every $\kk$-type $V$, see for instance Lemma 3.3 in [PS2]. Therefore Lemma 5.7 implies the first statement. The proof of the second statement is very similar to the proof of \refle{le00}.
\end{proof}

We now introduce the functors
\begin{eqnarray*}
\Ind: D^u-\mod &\hookrightarrow &A-\mod \\
M&\mapsto&A\otimes_{A_0}M,\\
\res_u:A-\mod&\hookrightarrow & D^u-\mod\\
M&\mapsto& D^u\otimes_{A_0}M.
\end{eqnarray*}
Obviously, $\res_u\circ\Ind=\id_{D^u-\mod}$.

\begin{lemma}\label{le95}
\[
\ker\gamma_u=\triplebrace{\Ann L_{u+1,1}=\Ann L_{-u-1,u+2}=\Ann L_{1,-u-2}}{\mathrm{~for~}u\notin \ZZ}{\Ann L_{-u-1,u+2}=\Ann L_{1,-u-2}}{\mathrm{~for~} u\in \ZZ_{\geq-1}}{\Ann L_{u+1,1}=\Ann L_{-u-1,u+2}}{\mathrm{~for~}u\in \ZZ_{\leq-2}}.
\]
\end{lemma}
\begin{proof}
First we prove that $\ker\gamma_u\subset\Ann L_{a,b}$ with $a,b$ as
in the statement.  Note that
$\res_u(t^u\CC[t^{\pm 1},x,y])$ contains a submodule generated by
$t^u$ isomorphic to $L_{u+1,1}$, $\res_u(x^u\CC[t^{\pm 1},x^{\pm
  1},y])/\res_u(x^u\CC[t,x^{\pm 1},y])$ contains a submodule with
highest vector $t^{-1}x^{u+1}$ isomorphic to $L_{-u-1,u+2}$ and
$\res_u(y^u\CC[t^{\pm 1},x^{\pm 1},y^{\pm 1}])/(\res_u(y^u\CC[t^{\pm
  1},x,y^{\pm 1}])+\res_u(y^u\CC[t,x^{\pm 1},y^{\pm 1}]))$ contains a
submodule with highest vector $t^{-1}x^{-1}y^{u+2}$ isomorphic to
$L_{1,-u-2}$. Hence $\ker\gamma_u \subset \ann L_{a,b}$. Next we see
from \refle{le1021sl} that all proper two-sided ideals of $D^u$ have
finite codimension. Thus, $\gamma_u (\Ann L_{a,b})$ is either $0$ or
has finite codimension in $D^u$. The latter is impossible because
$L_{a,b}$ is infinite-dimensional. Hence $\ker\gamma_u=\Ann L_{a,b}$.
\end{proof}

Since the eigenvalues of $\ad_H$ in $U(\gg)$ are all even, every
simple $(\gg,\kk)$-module is either odd or even.

As follows from \refle{le95}, all simple bounded $(\gg,\kk)$-modules with
central character $\chi(u+1,u)$ are $(D^u,\kk)$-modules. This allows us to first classify the simple $(D^u,\kk)$-modules and then use translation
functors to classify the bounded simple modules with arbitrary possible central character, see \refcor{corPrimIdeal}. 

Note that the functor $\ind$ maps $(D^u,\kk)$-mod into $\displaystyle \widetilde{(A, \kk)}$-mod, the latter being defined as the full subcategory of $A-\mod$ consisting of $\kk$-locally finite $A$-modules with semisimple action of $\EE$.
\begin{lemma}\label{le89}
For any simple $(D^u,\kk)$-module $M$ there exists a simple $\widetilde{(A, \kk)}$-module $\hat M$ with $\res_u(\hat M)\simeq M$.
\end{lemma}
\begin{proof}
Let $N$ be a maximal proper $A$-submodule of $\ind(M)$. Then $\res_u(N)\ncong M$ as $M$ generates $\ind(M)$. Therefore $\res_u(N)=0$ and one defines $\hat M$ as $\ind(M)/N$.
\end{proof}

Set $f:=x^2-2ty$, $\Delta:=\partial_x^2-2\partial_y\partial_t$ and note that $f,\Delta\in A^\kk$. For every fixed $p\in\CC$,we put $R^p:=f^p\CC[t,x,y,f^{-1}]$. Then clearly $R^p$ is an $(A,\tilde\kk)$-module and $\res_u (R^p)=0$ if $u-2p\notin\ZZ$. Otherwise,
\begin{equation}\label{eqle89}
\res_u(R^p)=\doublebrace{\CC f^{\frac{u}{2}}\oplus f^{\frac{u-2}{2}}\HH_2\oplus f^{\frac{u-4}{2}}\HH_4\oplus\dots}{\text{~for~} u-2p\in2\ZZ}{\CC f^{\frac{u-1}{2}}\HH_1\oplus f^{\frac{u-3}{2}}\HH_3\oplus f^{\frac{u-5}{2}}\HH_4\oplus\dots}{\text{~for~} u-2p\in2\ZZ+1,}
\end{equation}
where $\HH_n$ denotes the space of homogeneous polynomials of degree $n$ in $\CC[t,x,y]$ annihilated by $\Delta$ (as a $\kk$-module $\HH_n$ is isomorphic to $V_{2n}$).
\begin{lemma}\label{le910}
\item[(a)] For $u\notin\ZZ$ and for $u=-1,-2$, $\res_u(R^{\frac{u}{2}})$ and $\res_u(R^{\frac{u+1}{2}})$ are simple $D^u$-modules.
\item[(b)] For $u\in 2\ZZ_{\geq 0}$, $\res_u(R^{\frac{u+1}{2}})$ is a simple $D^u$-module and there is an exact sequence
\begin{equation}\label{eqle910+even}
0\to V_{u,0}\to\res_u(R^{\frac{u}{2}})\to I^{+}_{u,0}\to 0
\end{equation}
for some simple $D^u$-module $I^+_{u,0}$.
\item[(c)] For $u\in 1+2\ZZ_{\geq 0}$, $\res_u(R^{\frac{u}{2}})$ is a simple $D^u$-module and there is an exact sequence
\begin{equation*}\label{eqle910+odd}
0\to V_{u,0}\to\res_u(R^{\frac{u+1}{2}})\to I^{-}_{u,0}\to 0
\end{equation*}
for some simple $D^u$-module $I^-_{u,0}$.
\item[(d)] For $u\in 2\ZZ_{\leq -2}$, $\res_u(R^{\frac{u}{2}})$ is a simple $D^u$-module and there is an exact sequence
\begin{equation*}\label{eqle910-even}
0\to I^-_{u,0}\to\res_u(R^{\frac{u+1}{2}})\to V_{0,-3-u}\to 0
\end{equation*}
for some simple $D^u$-module $I^-_{u,0}$.
\item[(e)] For $u\in 1+2\ZZ_{\leq -1}$, $\res_u(R^{\frac{u+1}{2}})$ is a simple $D^u$-module and there is an exact sequence
\begin{equation*}\label{eqle910-odd}
0\to I^+_{u,0}\to\res_u(R^{\frac{u}{2}})\to V_{0,-3-u}\to 0
\end{equation*}
for some simple $D^u$-module $I^+_{u,0}$.
\end{lemma}
\begin{proof}
The isomorphism \refeq{eqle89} yields
\begin{equation}\label{eqle910res}
c(\res_u (R^{\frac{u}{2}}))=1+z^4+z^8+\dots,\quad c(\res_u (R^{\frac{u+1}{2}}))=z^2+z^6+z^{10}+\dots\quad.
\end{equation}
Thus, if $\res_u(R^{\frac{u}{2}})$ (respectively $\res_u(R^{\frac{u+1}{2}})$) is not simple it has a unique simple finite-dimensional submodule or a unique simple finite-dimensional quotient. By \refle{le95} the latter can happen only if $u\in \ZZ_{\geq 0}$ or $u\in \ZZ_{\leq-3}$. Hence (a).

Let $u\in 2\ZZ_{\geq 0}$. Then $\res_u(R^{\frac{u}{2}})$ contains $\res_u(\CC[t,x,y])\simeq V_{u,0}$ as a finite-dimensional simple submodule, hence \refeq{eqle910+even}. The $\gg$-module $\res_u(R^{\frac{u+1}{2}})$ has the same central character as $\res_u(R^{\frac{u}{2}})$ and, since $V_{n,0}$ is not a subquotient of $\res_u(R^{\frac{u+1}{2}})$ by \refeq{eqle910res}, $\res_u(R^{\frac{u+1}{2}})$ is a simple $D^u$-module. Hence (b).

As $\Delta(f^{-\half})=0$, $f^{-\half}$ generates a proper $A$-submodule $M\subset f^{\half}\CC[t,x,y,f^{-1}]$. A direct computation shows that $\dim\res_u(M)=\infty$ for any $u\in1+2\ZZ_{\geq -2}$. Furthermore, the only finite-dimensional module, whose central character coincides with that of $D^u$ is $V_{0,-3-u}$. Therefore one necessarily has
\[
0\to I_{u,0}^+\to\res_u(R^{\frac{u}{2}})\to V_{0,-3-u}\to 0
\]
where $I^+_{u,0}:=\res_u(M)$. $\res_u(R^{\frac{u+1}{2}})$ is simple by the same reason as in (b). Hence (e).

(c) and (d) are similar to (b) and (e).
\end{proof}

For any $u\in\CC$ we define now $I^+_{u,0}$ (respectively, $I^-_{u,0}$) as the unique simple infinite-dimensional constituent of $\res_u(R^{\frac{u}{2}})$ (resp., $\res_u(R^{\frac{u+1}{2}})$).

\begin{corollary}\label{cor103sl}
Every simple even infinite-dimensional $(D^u,\kk)$-module is isomorphic to $I^\pm_{u,0}$.
\end{corollary}
\begin{proof}
For every fixed $u$ and any sufficiently large $m\in 2\ZZ_{\geq 0}$ (such that $V_m$ is not a $\kk$-type of $V_{u,0}$ or $V_{0,-3-u}$ for $u\in\ZZ$), \refle{le910} implies $\hom_\kk(V_m,I^\pm_u)\neq 0$. The statement follows now from \refcor{cor98}.
\end{proof}

\begin{lemma}\label{le104sl}
If $u\notin \frac {1}{2}+\ZZ$, then every $(D^u,\kk)$-module is even.
\end{lemma}
\begin{proof}
Assume that $M$ is an odd simple $(D^u,\kk)$-module and $u\notin \frac {1}{2}+\ZZ $. Let $\hat M$ be as in \refle{le89}, $A_f$ denote the localization of $A$ in $f$, $\hat{M}_f=A_f\otimes_A \hat{M}$. First, we claim that if  $u\notin \frac
{1}{2}+\ZZ$, then $\hat{M}_f\neq 0$. Indeed, $\hat{M}_f=0$ implies
that $f$ acts
locally nilpotently on $\hat{M}$.
Then $M^0:=\Ker f$ is a $\kk$-submodule of $\hat{M}$ and a
straightforward calculation using \refeq{princ} shows
$\Omega_{| M^0}=2(\EE+3)(\EE+2)_{| M^0}$. Thus $\hom_{\kk}(V_m,M^0)\neq 0$ only if $2(d+3)(d+2)=\frac{m^2}{2}+m$ or equivalently $(d+\frac{5}{2})^2=(\frac{m+1}{2})^2$, where $d$ is the eigenvalue of $\EE$ on $M^0$. Since $d\in u+\ZZ$, $u\notin \frac{1}{2}+\ZZ $ implies $M^0=0$.

Our next observation is that $\hat{M}_f$ is an odd $(A,\kk)$-module
and that $t$ does not act locally nilpotently on $\hat{M}_f$. Indeed, if $t$ acts locally nilpotently, by $\kk$-invariance $x$ and $y$ act locally nilpotenly, and therefore $f$ acts locally
nilpotently. Contradiction. Therefore  $\hat{M}_f$ is a submodule of its localization in $t$, $\hat{M}_{f,t}$. Furthermore, for some odd $m$ there
exists a non-zero vector $v\in \hat{M}_{f,t}$ such that $H\cdot v=m
v$, $E\cdot v=0$
and  $\EE\cdot v=uv$. The expressions for $E,H$ and $\EE$ imply
\[
\partial_t v=\frac{-(u+m/2)ty+mx^2/2}{tf} v, \partial_x v=\frac{(u-m/2)x}{f} v, \partial_y v=\frac{(m/2-u)t}{f} v.
\]
Thus, every vector in $\hat{M}_{f,t}$ can be obtained from $v$ by 
applying elements of $\CC[t^{\pm 1},x,y,f^{-1}]$, i.e.
$\hat{M}_{f,t}=\CC[t^{\pm 1},x,y,f^{-1}]v$. It is not difficult to see
that $v=t^{\frac{m}{2}}f^{\frac{2u-m}{4}}$ satisfies the above
relations. The $A_{f,t}$-module 
$\CC[t^{\pm 1},x,y,f^{-1}]v$ is simple and free over
$\CC[t^{\pm 1},x,y,f^{-1}]$. Hence
$\hat{M}_{f,t}\simeq \CC[t^{\pm 1},x,y,f^{-1}]v$ and it is
obvious that $\hat{M}_{f,t}$ has no non-zero $\kk$-finite
vectors. As we pointed out above, $\hat{M}_f\subset
\hat{M}_{f,t}$. Therefore $\hat{M}_f=0$.
\end{proof}

We now turn to odd simple $(D^u,\kk)$-modules.

\begin{lemma}\label{le105odd} Let $u\in \frac {1}{2}+\ZZ$. Up to isomorphism, there exists exactly one odd simple $(D^u,\kk)$-module $J_{u,0}$. Moreover,
\begin{equation}\label{eqcJu0}
c(J_{u,0})=\doublebrace{z^{2-2u}+z^{6-2u}+z^{10-2u}+\dots}{\text{~for~}u<0}{z^{4+2u}+z^{8+2u}+z^{12+2u}+\dots}{\text{~for~}u>0}.
\end{equation}

\end{lemma}
\begin{proof} Let $P\subset G=SL(3)$ be the maximal parabolic subgroup whose Lie algebra $\pp$ equals $\bb\oplus\gg^{-\alpha_1}$, $K\subset G$ be the algebraic subgroup with Lie algebra $\kk$, and $Z$ be the closed $K$-orbit on $G/P\simeq\PP^2$. Then $Z\simeq \PP^1$ and the embedding $i:Z \to \PP^2$ is a Veronese embedding of degree 2. It is not difficult to verify that the relative tangent bundle $\mathcal{T}_P$ of the projection $p:G/B\to G/P$ is a $\OO_{G/B}$-submodule of the twisted sheaf of differential operators $\DD_{G/B}^{(u+1)\omega_1+\omega_2}$(the definition of $\DD_{G/B}^{\zeta}$ see in [PS2], Section 5). Furthermore, the direct image $p_*(\DD_{G/B}^{(u+1)\omega_1+\omega_2}/\mathcal{I}_P)$, where $\mathcal{I}_P$ is the left ideal in $\DD_{G/B}^{(u+1)\omega_1+\omega_2}$ generated by $\mathcal{T}_P$, is a well-defined twisted sheaf of differential operators on $G/P$. We denote this sheaf by $\DD_{G/P}^{(u+1)\omega_1+\omega_2}$.

Our next observation is that, similarly to the equivalence of categories $i_\bigstar$ discussed in Section 5 of [PS2], Kashiwara's theorem yields an equivalence of categories
\[
i_\bigstar^u:\OO_Z(2u)\otimes_{\OO_{G/P}}\DD_{G/P}\otimes_{\OO_{G/P}} \OO_Z(-2u)~-~\mod\to (\DD_{G/P}^{(u+1)\omega_1+\omega_2}~-~\mod)^Z,
\]
where $(\DD_{G/P}^{(u+1)\omega_1+\omega_2}~-~\mod)^Z$ denotes the full subcategory of $\DD_{G/P}^{(u+1)\omega_1+\omega_2}$- mod supported on $Z$, and $\OO_Z(2u)$ is the line bundle on $Z$ with Chern class $2u$. Therefore we can put
\[
J_{u,0}:=\Gamma(\PP^2,i_\bigstar^u\OO_Z(2u)).
\]
It is clear that $J_{u,0}$ is a $(\gg,\kk)$-module, and furthermore, using the facts that $\mathcal N\simeq\OO_Z(4)$ and that $i_\bigstar^u\OO_Z(2u)$ has a filtration with successive quotients $\OO_Z(2u+4(i+1))$, one easily verifies that $c({J_{u,0}})$ is given by the right-hand side of \refeq{eqcJu0}. Since there are no finite-dimensional modules with central character $\chi(u+1,1)$ for $u\in\half+\ZZ$, $J_{u,0}$ is a simple $\gg$-module.

It remains to prove that every simple odd $(D^u,\kk)$-module is
isomorphic to $J_{u,0}$ for some $u\in\half+\ZZ$. Let $M$ be a simple
odd $(D^u,\kk)$-module and $\hat M$ be a simple $(A,\tilde\kk)$-module
such that $\res_u (\hat M)=M$. Then by the proof of \refle{le105odd}
$\hat{M}_f=0$. For every $\bb_\kk$-highest vector $v\in \res_u(\hat
M)$ there exists $k$ such that $f^k\cdot v=0$. Let $v$ have weight
$m$. Then by the relation $(d+\frac{5}{2})^2=(\frac{m+1}{2})^2$ from
the proof of \refle{le105odd}, $\frac{m+1}{2}=\pm(u+2k+\frac{5}{2})$,  
as $\EE f^k\cdot v=(2k+u)f^k\cdot v$. Without loss of generality we
may assume that $m$ is very large and then 
$\frac{m+1}{2}=(u+2k+\frac{5}{2})$. 
Therefore $\Hom_\kk(V_m, M)\neq 0$ implies $m=2u+4k+4$. Hence if $M$ and $M'$ are two odd $(D^u,\kk)$-modules one can find $m$ such that $\Hom_\kk(V_m,M)\neq 0$, $\Hom_\kk(V_m, M')\neq 0$. But then $M\simeq M'$ by \refcor{cor98}.
\end{proof}

Let $M$ be some $A$-module with semisimple $\EE$-action. Consider the 
$U(\gg)$-modules $M^{(n)}:=M\otimes S^n(\span\{x,y,t\})$ for
$n\in\mathbb Z_{\geq 0}$, together with the linear operators
\begin{eqnarray*}
\bar d: M^{(n)}&\to &M^{(n-1)}\\
\bar d&=&t\otimes\partial_t+x\otimes\partial_x+y\otimes\partial_y\\
\bar \delta: M^{(n)}&\to &M^{(n+1)}\\
\bar\delta&=&\partial_t\otimes t+\partial_x\otimes x+\partial_y\otimes y.
\end{eqnarray*}
It is straightforward to check that $\bar d$, $\EE\otimes 1-1\otimes \EE$
and $\bar\delta$ form a standard $\sl(2)$-triple. Let $\res_s(M^{(k)})$ be
the eigenspace of the operator $\EE\otimes 1+1\otimes \EE$ in
$M^{(k)}$. Then obviously $\bar d$ and $\bar \delta$ induce operators
\begin{eqnarray*}
d: \res_s (M^{(n)})&\to &\res_s( M^{(n-1)})\\
\delta: \res_s (M^{(n-1)})&\to &\res_s( M^{(n)}),
\end{eqnarray*}
and elementary $\sl(2)$ representation theory implies 
that if $s\notin\ZZ$, $s<n-1$ or $s\geq 2n$, then $d$ is surjective, $\delta$ is injective, and
\begin{equation}\label{eq87}
\res_s( M^{(n)})=\Ker d\oplus \im \delta.
\end{equation}
For any $(D^u,\kk)$-module $M$ choose a simple $(A,\tilde\kk)$-module $\hat M$ such that $\res_u(\hat M)=M$ (in fact $\hat M$ is unique).

Let $T^n(M):=\res_{u+n}({\hat M}^{(n)})\cap \Ker d$. If $u\neq
-1,0,\dots,n-1$, \refeq{eq87} implies
\begin{equation}\label{eqcTnM}
c(T^n(M))=c(\res_{u+n}(\hat M^{(n)}))-c(\res_{u+n}(\hat M^{(n-1)})).
\end{equation}

\begin{lemma}\label{le811}
Let $M$ be a bounded simple $(D^u,\kk)$-module. Assume that $u\neq -1,0,\dots,n-1$. Then $T^n(M)$ is a simple $(\gg,\kk)$-module with central character $\chi(u+1-n,n+1)$.
\end{lemma}
\begin{proof}
\refle{le95} implies that $M$ is a $(\gg,\kk)$-module with central character $\chi(u+1,1)$. Therefore $M\otimes S^n(\span\{x,y,t\})$ has constituents with central character $\chi(u+1+n-2k,1+k)$, $k=0,\dots,n$, and $\im \delta$ has constituents with central character $\chi(u+1+n-2k,1+k)$, $k=0,\dots,n-1$. Thus, $T^n(M)$ is a direct summand of $M\otimes S^n(\span\{x,y,t\})$ with central character $\chi(u+1-n,n+1)$.

Our restrictions on $u$ imply that the weights $(u+1)\omega_1+\omega_2$ and $(u-n+1)\omega_1+(n+1)\omega_2$ belong to the same Weyl chamber and have the same stabilizer in the Weyl group. Hence, $T^n$ is nothing but the translation functor
\[
T_{(u+1)\omega_1+\omega_2}^{(u-n+1)\omega_1+(n+1)\omega_2}:\BB_\kk^{\chi(u+1,1)}\to\BB_\kk^{\chi(u-n+1,n+1)}.
\]
 

Therefore $T^n$ is an equivalence of categories, in particular $T^n(M)$ is simple.
\end{proof}

We put for $u\neq -1,0,\dots,n-1$
\[
I^\pm_{u,n}:=T^n(I^\pm_{u,0}),
\]
\[
J_{u,n}:=T^n(J_{u,0}).
\]

\begin{theorem}\label{th914}
Let $M$ be a simple bounded infinite-dimensional $(\gg,\kk)$-module with central character $\chi$. Then
\item[(a)] if $\chi=\chi(u+1-n,n+1)$ for $u\notin\ZZ$,
\[
M\simeq\doublebrace{I^\pm_{u,n}}{\text{~for~}u\notin\half+\ZZ}{I^\pm_{u,n},J_{u,n}}{\text{~for~}u\in\half+\ZZ};
\]
\item[(b)] if $\chi=\chi(u+1-n,n+1)$ for $u\in\ZZ_{\geq n}$,
\[
M\simeq I^\pm_{-n-3,u-n},I^\pm_{u,n};
\]
\item[(c)] if $\chi=\chi(-1-n,n+1)$,
\[
M\simeq I^\pm_{-2,n};
\]
\item[(d)] if $\chi=\chi(0,n+1)$,
\[
M\simeq(I^\pm_{-2,n})^\tau,
\]
where $\tau$ stands for the outer automorphism $\tau(X)=-X^t$ for any 
$X\in\gg$.
\end{theorem}
\begin{proof}
By \refcor{corPrimIdeal} every simple bounded $(\gg,\kk)$-module has
central character $\chi$ of the form $\chi(u+1-n,n+1)$ for some
$n\in\ZZ_{\geq 0}$ and some $u\in\{\CC\backslash
\ZZ_{<n-1}\}\cup\{-2\}$. Moreover,
$T^n=T_{(u+1)\omega_1+\omega_2}^{(u-n+1)\omega_1+(n+1)\omega_2}$ is an
equivalence of the categories $\BB_\kk^{\chi(u+1,1)}$ and
$\BB_\kk^{\chi(u+1-n,n+1)}$. If $u\notin\ZZ$, $\half+\ZZ$ then
$\BB_\kk^{\chi(u+1,1)}$ has two non-isomorphic simple objects, and, if $u\in \half+\ZZ$, $\BB_\kk^{\chi(u+1,1)}$ has three non-isomorphic simple objects. This implies (a).

If $u\in\ZZ_{\geq 0}$, $u\geq n$, we have $\chi=\chi(u+1-n,n+1)=\chi((-n-3)+1-(u-n),(u-n)+1)$, hence in this case $\BB_\kk^\chi$ has 4 non-isomorphic simple objects: $I^\pm_{u,n}$ and $I^\pm_{-n-3,u-n}$. This proves (b). If $n=-2$, $\BB_\kk^\chi$ is equivalent to $\BB_\kk^{\chi(1,1)}$ and has two simple objects, $I^\pm_{-2,n}$, which proves (c). Finally if $u=n-1$, the automorphism $\tau$ establishes an equivalence between $\BB_\kk^{\chi(0,n+1)}$ and $\BB_\kk^{\chi(-1-n,n+1)}$, hence (d).
\end{proof}

\begin{lemma}\label{le915}
For $a\in\ZZ_{\geq 2}$, define
\[
\mu_n(a,z):=\frac{z^a}{1-z^4}\otimes c(V_{n,0})-\frac{z^{a-2}}{1-z^4}\otimes c(V_{n-1,0}).
\]
For $a\in\ZZ_{\geq 0}$, define
\[
\kappa_n(a,z):=\frac{z^a}{1-z^4}\otimes c(V_{n,0})-\frac{z^{a+2}}{1-z^4}\otimes c(V_{n-1,0}).
\]
Then
\begin{equation}\label{eq916x}
\mu_{2p}(a,z)=\frac{z^a}{1-z^4}+\frac{z^{a-2}(z^4+z^8+\dots+z^{4p})}{1-z^2},
\end{equation}
\begin{equation}\label{eq916x+1}
\mu_{2p+1}(a,z)=\frac{z^{a}(1+z^4+\dots+z^{4p})}{1-z^2},
\end{equation}
\begin{equation}\label{eq916x+2}
\kappa_{2p}(a,z)=\frac{z^a}{1-z^4}+\frac{z^{|a-4|}+\dots+z^{|a-4p|}}{1-z^2},
\end{equation}
\begin{equation}\label{eq916x+3}
\kappa_{2p+1}(a,z)=\frac{z^{|a-2|}+\dots+z^{|a-4p-2|}}{1-z^2}.
\end{equation}
\end{lemma}
\begin{proof}
Since $V_{n,0}=S^n(V_{1,0})$, and since $S^n(V_{1,0})$ is isomorphic as a $\kk$-module to $S^n(V_2)$, we have
\[
c(V_{2p,0})=1+z^4+\dots+z^{2p},
\]
\[
c(V_{2p+1,0})=z^2+z^6+\dots+z^{2p+2}.
\]
Recall that $z^a\otimes z^b=\pi(z^a\sum_{i=0}^{i=b}z^{b-2i})$ (\refle{le202},(b)). Therefore
\begin{eqnarray*}
\frac{z^a}{1-z^4}\otimes z^{2k}-\frac{z^{a-2}}{1-z^4}\otimes z^{2k-2}&=&\pi\left(\frac{z^{a-2}(z^2\sum\limits_{i=0}^{i=2k}z^{2k-2i}-z^{-2}\sum\limits_{i=0}^{i=2k-2}z^{2k-2i})}{1-z^4}\right)=\\
&=&\pi\left(\frac{z^{a-2}(z^{2k+2}+z^{2k})}{1-z^4}\right)=\frac{z^{a-2+2k}}{1-z^2}.
\end{eqnarray*}
\begin{eqnarray*}
\frac{z^a}{1-z^4}\otimes z^{2k}-\frac{z^{a+2}}{1-z^4}\otimes z^{2k-2}&=&\pi\left(\frac{z^{a}(\sum\limits_{i=0}^{i=2k}z^{2k-2i}-z^2\sum\limits_{i=0}^{i=2k-2}z^{2k-2-2i})}{1-z^4}\right)=\\
&=&\pi\left(\frac{z^{a}(z^{-2k}+z^{2-2k})}{1-z^4}\right)=\pi\left(\frac{z^{a-2k}}{1-z^2}\right)=\frac{z^{|a-2k|}}{1-z^2}.
\end{eqnarray*}
The above identities imply \refeq{eq916x}-\refeq{eq916x+3}.
\end{proof}
%\begin{eqnarray*}
%\mu_{2p}(a,z)&=&\frac{z^a}{1-z^4}\otimes(1+z^4+\dots+z^{4p})-\\
%&&-\frac{z^{a-2}}{1-z^4}\otimes(z^2+\dots+z^{4p-2})=\\
%&=&\frac{z^a}{1-z^4}+\frac{z^{a-2}(z^4+z^8+\dots+z^{4p})}{1-z^2}
%\end{eqnarray*}
%\begin{eqnarray*}
%\mu_{2p+1}(a,z)&=&\frac{z^a}{1-z^4}\otimes(z^2+z^6+\dots+z^{4p+2})-\\
%&&-\frac{z^{a-2}}{1-z^4}\otimes(1+z^4+\dots+z^{4p})=\\
%&=&\frac{z^{a}(1+z^4+\dots+z^{4p})}{1-z^2}
%\end{eqnarray*}
%\begin{eqnarray*}
%\kappa_{2p}(a,z)&=&\frac{z^a}{1-z^4}\otimes(1+z^4+\dots+z^{4p})-\\
%&&-\frac{z^{a+2}}{1-z^4}\otimes(z^2+\dots+z^{4p-2})=\\
%&=&\pi\left(\frac{z^{a}(z^{-4}+\dots+z^{-4p})}{1-z^2}\right)+\frac{z^a}{1-z^4}=\\
%&=&\frac{z^{|a-4|}+\dots+z^{|a-4p|}}{1-z^2}+\frac{z^a}{1-z^4}
%\end{eqnarray*}
%\begin{eqnarray*}
%\kappa_{2p+1}(a,z)&=&\frac{z^a}{1-z^4}\otimes(z^2+z^6+\dots+z^{4p+2})-\\
%&&-\frac{z^{a+2}}{1-z^4}\otimes(1+z^4+\dots+z^{4p})=\\
%&=&\pi\left(\frac{z^{a}(z^{-2}+z^{-6}+\dots+z^{-4p-2})}{1-z^2}\right)=\\
%&=&\frac{z^{|a-2|}+\dots+z^{|a-4p-2|}}{1-z^2}
%\end{eqnarray*}

\begin{theorem}\label{th916}
\item[(a)]Let $u\notin \ZZ$, $\half+\ZZ$. Then
\[\begin{array}{l}
c(I^+_{u,n})=\kappa_n(0,z), \quad c(I^-_{u,n})=\mu_{n}(2,z).
\end{array}\]
\item[(b)]Let $u\in \half+\ZZ$. Then
\[\begin{array}{ll}
c(J_{u,n})=\kappa_n(4+2u,z)&\text{~for~}u\geq -\half,\\
c(J_{u,n})=\mu_n(2-2u,z)&\text{~for~}u\leq -\threehalfs.
\end{array}\]
\item[(c)]Let $u\in 2\ZZ_{\geq 0}$. Then
\[\begin{array}{ll}
c(I^+_{u,0})=\frac{z^{2u+4}}{1-z^4},&c(I^-_{u,0})=\frac{z^{2}}{1-z^4},\\
c(I^+_{u,n})=\kappa_n(2u+4,z),&c(I^-_{u,n})=\mu_n(2,z).
\end{array}\]
\item[(d)]Let $u\in 1+2\ZZ_{\geq 0}$. Then
\[\begin{array}{ll}
c(I^+_{u,0})=\frac{1}{1-z^4},&c(I^-_{u,0})=\frac{z^{2u+4}}{1-z^4},\\
c(I^+_{u,n})=\kappa_n(0,z),&c(I^-_{u,n})=\kappa_n(2u+4,z).
\end{array}\]
\item[(e)]Let $u\in 2\ZZ_{\leq-2}$. Then
\[\begin{array}{ll}
c(I^+_{u,0})=\frac{1}{1-z^4},&c(I^-_{u,0})=\frac{z^{-2-2u}}{1-z^4},\\
c(I^+_{u,n})=\kappa_n(0,z),&c(I^-_{u,n})=\mu_n(-2-2u,z).
\end{array}\]
\item[(f)]Let $u\in -1+2\ZZ_{\leq-1}$. Then
\[\begin{array}{ll}
c(I^+_{c,0})=\frac{z^{-2-2u}}{1-z^4},&c(I^-_{u,0})=\frac{z^{2}}{1-z^4},\\
c(I^+_{u,n})=\mu_n(-2-2u,z),&c(I^-_{u,n})=\mu_n(2,z).
\end{array}\]
\item[(g)]
\[\begin{array}{l}
c(I^+_{-2,n})=c((I^+_{-2,n})^\tau)=\kappa_n(0,z),\\
c(I^-_{-2,n})=c((I^-_{-2,n})^\tau)=\mu_n(2,z).
\end{array}\]
\end{theorem}
\begin{proof}
Using \refeq{eqcTnM} one obtains the identities
\begin{equation}\label{eqth916}
\begin{array}{l}
c(I^\pm_{u,n})=c(I^\pm_{u,0}\otimes V_{n,0})-c(I^\mp_{u+1,0}\otimes V_{n-1,0}),\\
c(J_{u,n})=c(J_{u,0}\otimes V_{n,0})-c(J_{u+1,0}\otimes V_{n-1,0}).
\end{array}
\end{equation}
The theorem is a straightforward corollary of \refeq{eqth916}. Indeed, let us prove (f). In this case
\[
c(I_{u,0}^+)=\frac{z^{-2u-2}}{1-z^4},\quad c(I_{u-1,0}^+)=\frac{z^{-2u-4}}{1-z^4},
\]
\[
c(I_{u,n}^+)=\frac{z^{-2u-2}}{1-z^4}\otimes c(V_{n,0})-\frac{z^{-2u-4}}{1-z^4}\otimes c(V_{n-1,0})=\mu_n(-2-2u,z);
\]
\[
c(I_{u-1,0}^-)=\frac{z^{-2u-4}}{1-z^4},\quad c(I_{u-1,0}^+)=\frac{1}{1-z^4},
\]
\[
c(I_{u,n}^-)=\frac{z^{2}}{1-z^4}\otimes c(V_{n,0})-\frac{1}{1-z^4}\otimes c(V_{n-1,0})=\mu_n(2,z).
\]
In all other cases the arguments are similar.
\end{proof}

\begin{corollary}
\item[(a)] The minimal $\kk$-type can be any $V_k$ but its multiplicity is always 1.
\item[(b)] For sufficiently large $i$ $c_i(M)=c_{i+4}(M)$ for any
  simple bounded $(\gg,\kk)$-module, and for sufficiently large $j$
  there are the following $\kk$-multiplicities:
\[
c_{4j}(I^\pm_{u,2p+1})=c_{4j+2}(I^\pm_{u,2p+1})=p+1,
\]
\[
c_{4j}(I^+_{u,2p})=p+1, c_{4j+2}(I^+_{u,2p})=p,
\]
\[
c_{4j+2}(I^-_{u,2p})=p+1, c_{4j}(I^-_{u,2p})=p,
\]
\[
c_{4j+1}(J_{u,2p+1})=c_{4j+3}(J_{u,2p+1})=p+1,
\]
\[
c_{4j+2u}(J_{u,2p})=p, c_{4j+2u+2}(J_{u,2p})=p+1.
\]
\item[(c)] The only multiplicity-free simple infinite-dimensional 
 $(\gg,\kk)$-modules are $I^\pm_{u,0}$, $J_{u,0}$, $I^\pm_{u,1}$,
 $J_{u,1}$, $(I^\pm_{-2,1})^\tau$.
\end{corollary}

The complete list of multiplicity-free simple $(\gg,\kk)$-modules has
been first found by Dj. Sijacki, see \cite{S} and the references
therein for a historic perspective on this problem. 

\section{Classification of simple bounded $(\sp(4),\sl(2))$-modules}\label{se6}
In this section we classify all simple bounded $(\gg,\kk)$-modules,
where $\gg=\sp(4)$ and $\kk$ is a principal $\sl(2)$-subalgebra or a
$\sl(2)$-subalgebra corresponding to a short root. We fix a Cartan
subalgebra $\hh\subset\gg$ and write the roots of $\gg$ as $\{\pm
2\eps_1,\pm 2\eps_2, \pm\eps_1\pm\eps_2\}$. Our fixed simple roots are
$\eps_1-\eps_2, 2\eps_2$, and $\rho=2\eps_1+\eps_2$. By $e_1$, $e_2$,
$h_1$, $h_2$, $f_1$, $f_2$ we denote the Serre generators of $\gg$
associated to our choice of simple roots, \cite{OV}. We define two
$\sl(2)$-subalgebras of $\gg$: one with basis $e_1$, $h_1$, $f_1$ and
one with basis $e_1+2e_2$, $3h_1+4h_2,3f_1+2f_2$. The first one is the
root subalgebra corresponding to the simple root $\eps_1-\eps_2$, and
the second one is a principal $\sl(2)$-subalgebra. In Sections
\ref{se6} and \ref{secKchars4-2}, we denote by $\kk$ any one of these
two subalgebras, referring respectively to the \emph{root case} and to
the \emph{principal case} when we want to be specific. We set
$\bb_\kk:=\bb\cap \kk$, where $\bb$ is the Borel subalgebra 
generated by $e_1$, $e_2$, $h_1$, $h_2$. By
$L_{a,b}$ we denote the simple $\bb$-highest weight $\gg$-module with
highest weight $a\eps_1+b\eps_2-\rho=(a-2)\eps_1+(b-1)\eps_2$, by
$V_{a,b}$ we denote the simple finite-dimensional $\gg$-module with
highest weight $a\eps_1+b\eps_2$, and $\chi(a,b)$ is the central character of $L_{a,b}$.

\begin{lemma}\label {le100}
Let dim $L_{a,b}=\infty$ and $\gkdim L_{a,b}\leq 2$. Then $a>| b|$ and $a,b\in \half+\ZZ$.
\end{lemma}
\begin {proof} Let $\lambda=a\eps_1+b\eps_2$. If $(\lambda,\alpha)\notin
  \ZZ_{>0}$ for all positive roots $\alpha$, then $L_{a,b}$ is
  a Verma module and therefore its Gelfand-Kirillov
  dimension equals 4. If $(\lambda,\check\alpha)\in \ZZ_{>0}$ for exactly one positive
  root, then one has the following exact sequence
\[
0\to L_{w_\alpha(\lambda)} \to M_{\lambda} \to L_{\lambda}\to 0,
\]
where $w_{\alpha}$ denotes the reflection in $\alpha$. A straightforward computation shows that in this case $\gkdim
L_{\lambda}=3$. Therefore $\gkdim L_{\lambda}\leq 2$ implies the existence
of two positive roots $\alpha$ and $\beta$ such that
$(\lambda,\check\alpha), (\lambda,\check\beta) \in \ZZ_{>0}$. One can see
immediately that at least one of these roots, say $\alpha$, is
simple. If $N_{\lambda}$ denotes the quotient of $M_{\lambda}$ by the 
submodule generated 
by a highest vector with weight $w_\alpha(\lambda)-\rho$, then $\gkdim
N_{\lambda}=3$. The condition $\gkdim L_{\lambda}\leq 2$ implies the 
reducibility of
$N_{\lambda}$ which in turn implies $(\lambda,\check\gamma) \in \ZZ_{>0}$
for the positive root $\gamma$ orthogonal to $\alpha$. That leaves only
two possibilities for $\lambda$: $\lambda$ is either regular integral
or $\lambda$ satisfies the conditions of the Lemma.

It remains to eliminate the case of a regular integral non-dominant $\lambda$. By using the translation functor we may assume without loss of generality that $\lambda$ belongs to the Weyl group orbit of $\rho$. That leaves four possibilities for $\lambda$: $2\eps_1-\eps_2$, $\eps_1-2\eps_2$, $\eps_1+2\eps_2$, $-\eps_1+2\eps_2$.  Let $\pp_1$ and $\pp_2$ be the parabolic subalgebras obtained from $\bb$ by joining $\eps_2-\eps_1$ and $-2\eps_2$ respectively. It is  not difficult to verify the existence of embeddings
\[
L_{2,-1}\to U(\gg)\otimes_{U(\pp_1)}F^1_{2,1},\quad
L_{1,-2}\to U(\gg)\otimes_{U(\pp_1)}F^1_{2,-1},
\]
\[
L_{1,2}\to U(\gg)\otimes_{U(\pp_2)}F^2_{2,1},\quad
L_{-1,2}\to U(\gg)\otimes_{U(\pp_2)}F^2_{1,2},
\]
where $F_{a,b}^1$ (respectively, $F_{a,b}^2$) is the finite
 dimensional $\pp_1$-module (resp., $\pp_2$-module) with $\bb$-highest
 weight $a\eps_1+b\eps_2-\rho$. Therefore 
 the Gelfand-Kirillov dimension of any of the above four simple
 modules equals the Gelfand-Kirillov dimension of the corresponding
 parabolically induced module, i.e. 3. The proof is now complete.
\end{proof}

\begin {corollary}\label{cor101}
Let $M$ be a simple bounded infinite-dimensional \gkm. Then $\ann M= \ann L_{a,b}$ for some $a,b$ with $a>|b|$, $a,b\in\half+\ZZ$. In particular, $\chi(a,b)$ is the central character of $M$.
\end{corollary}
\begin {proof}
By Duflo's theorem, $\ann M = \ann L_{a,b}$ for some $a,b$. It is known that $\half\dim X_{L_{a,b}}=\gkdim L_{a,b}$, thus $\gkdim M\geq \gkdim L_{a,b}$. On the other hand, \mbox{$\displaystyle\gkdim M\leq 2=b_\kk$} holds by Theorem 4.4 in [PS2] . Hence $\gkdim L_{a,b}\leq 2$, and \refle{le100} applies to $L_{a,b}$.
\end{proof}

\begin {corollary}\label{cor102}
Let $a,b\in \half+\ZZ$, $a>|b|$. Then $\BB_\kk^{\chi(a,b)}$ is equivalent to $\BB_\kk^{\chi(\threehalfs,\half)}$.
\end {corollary}
\begin{proof}
It is well known that the categories $U^{\chi(a,b)}$-mod for $(a,b)$ as above are translation-equivalent to the category $U^{\chi^(\frac{3}{2},\frac{1}{2})}$-mod. Since the translation functor preserves the subcategories of bounded modules, the categories $\BB_\kk^{\chi(a,b)}$ and $\BB_\kk^{\chi(\threehalfs,\half)}$ are equivalent as well. 
\end{proof}

Our next step is to describe the quotient algebra $\univ(\gg) / \ann L_{\threehalfs,\half}$. In this section we denote by $A$ the Weyl algebra in two variables, i.e. the algebra of differential operators acting in $\CC[x,y]$. We introduce a $\ZZ_2$-grading, $A\eqdef A_0\oplus A_1$, by putting $\deg x=\deg y=\deg\del_x=\deg\del_y\eqdef \overline{1}\in\ZZ_2$. It is well known that there exists a surjective algebra homomorphism
\[
\kappa:\univ(\gg)\to A_0
\]
such that
\[
\kappa(e_1)=x\del_y,\quad \kappa (e_2)=\frac{y^2}{2},\quad \kappa (f_1)=y\del_x, \quad \kappa (f_2)=-\frac{\del_y^2}{2},
\]
\[
\kappa(h_1)=x\del_x-y\del_y,\quad \kappa(h_2)=y\del_y+\half.
\]
The kernel of $\kappa$ equals $\ann L_{\threehalfs,\half}$. Furthermore, $\kappa (\kk)$ is spanned by $E\eqdef x\del_y$, $F\eqdef y\del_x$, $H\eqdef x\del_x-y\del_y$ in the root case, and respectively by $E\eqdef x\del_y+y^2$, $H\eqdef 3x\del_x+y\del_y+2$, $F\eqdef 3y\del_x-\del_y^2$ in the principal case.

The problem of describing all simple modules in $\BB_\kk^{\chi(\threehalfs,\half)}$ is equivalent to the problem of describing all simple $(A_0,\kk)$-modules, i.e. all simple locally $\kappa(\kk)$-finite $A_0$-modules. The following lemma reduces this problem to a classification of all simple $(A,\kk)$-modules.

\begin{lemma}\label{le103}
Every simple $(A,\kk)$-module $M$ is a $\ZZ_2$-graded $A$-module, i. e. $M=M_0\oplus M_1$ where $M_0$ and $M_1$ are simple $(A_0,\kk)$-modules. Furthermore, $M=A\otimes_{A_0}M_0$, and the $\ZZ_2$-grading on $M$ is unique up to interchanging $M_0$ with $M_1$.
\end{lemma}
\begin{proof}
The element $H$ (as defined above separately for the root case and for the principal case) acts semisimply on $M$ with integer eigenvalues. We define $M_0$ (respectively, $M_1$) as the direct sum of $H$-eigenspaces with even (resp., odd) eigenvalues. It is obvious that $M=M_0\oplus M_1$, that $M_0$ and $M_1$ are simple $A_0$ modules, and that $M=A\otimes_{A_0}M_0$. Since $M_0$ and $M_1$ are non-isomorphic as $A_0$-modules, the uniqueness follows from the fact that a decomposition of $M$ as an $A_0$-module into a direct sum of two non-isomorphic $A_0$-modules is unique.
\end{proof}

\textbf{Remark.} More generally, if $\kk'$ is a subalgebra of $\gg'=\sp(2m)$ such that the centralizer of $\kk'$ in the Weyl $A'$ algebra of $m$ indeterminates is abelian, every $(A',\kk')$-module is a multiplicity-free $(\gg',\kk')$-module whose primitive ideal is a Joseph ideal. F. Knop has classified all such subalgebras $\kk'$, \cite{Kn2}, which makes us optimistic that this idea can eventually lead to a classification of simple bounded $(\gg',\kk')$-modules.

\medskip

Let $\Fou:A\to A$ be the automorphism defined by
\[
\Fou(x)\eqdef \del_x,\quad \Fou(y)\eqdef\del_y,\quad\Fou(\del_x)\eqdef -x,\quad \Fou(\del_y)\eqdef -y
\]
If $M$ is an $A$-module, we denote by $M^\Fou$ the twist of $M$ by $\Fou$.
\begin{theorem} \label{th500}\myLabel{th500}\relax
In the root case, any simple $(A,\kk)$-module is isomorphic to $\CC[x,y]$ or $\CC[x,y]^\Fou$.
\end{theorem}
\begin{proof}
Let $M$ be a simple $(A,\kk)$-module. Then there exists $ 0\neq v\in M $ such that $E\cdot v=0$, i.e. $ x\partial_{y}\cdot v=0 $. Hence either $ x $ or $ \partial_{y} $
act locally nilpotently on $ M $.

Assume first that $ \partial_{y} $ acts locally nilpotently on $ M $. Then $ \partial_{x}\in\left[{\mathfrak k},\partial_{y}\right] $
also acts locally nilpotenly on $ M $. Let $ A^{+} $ be the abelian subalgebra in $A$
generated by $ \partial_{x},\partial_{y} $. One can find $ 0\neq w\in M $ such that $ A^{+}\cdot w=0 $, and hence
\begin{equation}
M\cong A\otimes_{ A^{+}}\CC\cong\CC\left[x,y\right].
\notag\end{equation}

If $ x $ acts locally nilpotently on $ M $, one considers $ M^{\Fou} $ and reduces to the previous case.\end{proof}

\begin{corollary} \label{cor35}\myLabel{cor35}\relax
In the root case, up to isomorphism, there are exactly four simple $(\gg,\kk)$-modules with central character $\chi(\threehalfs,\half)$. As $\kk$-modules two of these modules are isomorphic to
\begin{equation}
V_{0}\oplus V_{2}\oplus V_{4}\oplus\dots\quad,
\notag\end{equation}
and the other two are isomorphic to
\begin{equation}
V_{1}\oplus V_{3}\oplus V_{5}\oplus\dots\quad .
\notag\end{equation}
\end{corollary}

\begin{theorem} \label{th600}\myLabel{th600}\relax
In the principal case, up to isomorphism, there exist exactly two simple $(A,\kk)$-modules and they have the following $\kk$-module decompositions:
\begin{equation}
V_{0}\oplus V_{3}\oplus V_{6}\oplus V_{9}\oplus\text{\dots , }\quad V_{1}\oplus V_{4}\oplus V_{7}\oplus V_{10}\oplus\dots\quad .
\notag\end{equation}
\end{theorem}
\begin{proof}
Note that $ \kk $ is a maximal subalgebra of $ \gg $. Hence, every element
$ g\in{\mathfrak g}\backslash\kk $ acts freely on a simple $ \left( A,\kk\right) $-module $ M $. In particular, $ x^{2} $ acts
freely on $ M $, and therefore $ x $ acts freely on $ M $. Let $ A_{x} $ be the localization
of $A$ in $ x $, and $ M_{x}:={A}_{x}\otimes_{{A}}M $. Then $ M\subset M_{x} $. Fix $0\neq m\in M $ with $ E\cdot m=0 $ and
$ H\cdot m=\lambda m $ for a minimal $ \lambda\in{\mathbb Z}_{\geq0} $. Since $E=x\del_y+y^2$ and $H=3x\del_x+y\del_y+2$, we have
\[
%original version below
%\partial_{y}\cdot m=-\frac{y^{2}}{x}\cdot m\text{, }\partial_{x}\cdot m=\left(-\frac{y^{3}}{3x^{2}}+\frac{\lambda-2}{3x}\right)\cdot m.
%suggested by Alexey correction below
\partial_{y}\cdot m=-\frac{y^{2}}{x}\cdot m\text{, }\partial_{x}\cdot m=\left(\frac{y^{3}}{3x^{2}}+\frac{\lambda-2}{3x}\right)\cdot m.
\notag
\]
Therefore, $ M_{x}={\mathbb C}\left[x,x^{-1},y\right]\cdot m $. Set
\begin{equation}
u_{\lambda}\eqdef x^{\frac{\lambda-2}{3}}\exp{}\left({\frac{-y^{3}}{3x}}\right).
\notag
\end{equation}
Then it is easy to see that $ M_{x} $ is isomorphic to $ {\mathcal F}_{\lambda}\eqdef{\mathbb C}\left[x,x^{-1},y\right]u_{\lambda} $ and that
$ {\mathcal F}_{\lambda}={\mathcal F}_{\lambda+3} $. Hence, $ M_{x} $ is isomorphic $ {\mathcal F}_{0},{\mathcal F}_{1} $ or $ {\mathcal F}_{2} $.

Next we calculate $ \Gamma_\kk\left({\mathcal F}_{\lambda}\right)
$. Note that the space of $ {\mathfrak b}_{{\mathfrak k}} $-singular
vectors in $ {\mathcal F}_{\lambda} $ is spanned by the family $
u_{\lambda+3k} $, $ k\in{\mathbb Z} $ of solutions to the
differential equation
\begin{equation}
E\cdot u=x\partial_{y}(u)+y^{2}u=0.
\notag
\end{equation}
If $ \lambda\in{\mathbb Z}_{\geq0} $, then $ F^{\lambda+1}\cdot u_{\lambda} $ is again a $ {\mathfrak b}_{{\mathfrak k}} $-highest vector of
weight $ -\lambda-2 $. Therefore $ F^{\lambda+1}\cdot u_{\lambda}=c u_{-\lambda-2} $ for some constant $ c $. On the other hand,
$ u_{-\lambda-2}\in{\mathcal F}_{\lambda} $ iff \mbox{$ \lambda-\left(-\lambda-2\right)=2\lambda+2\in3{\mathbb Z} $} or $ \lambda=3k+2 $. Hence $ F^{\lambda+1} \cdot u_{\lambda}=0 $ for $ \lambda=3k $ or $\lambda=3k+1 $.
Thus, $ \Gamma_\kk\left({\mathcal F}_{0}\right) $ is generated by $ u_{3k} $ for $ k\geq0 $,  $ \Gamma_\kk\left({\mathcal F}_{1}\right) $ is generated by $ u_{3k+1} $ for
$ k\geq0 $, and we have the $\kk$-module decompositions
\begin{equation}
\Gamma_\kk\left({\mathcal F}_{0}\right)\simeq V_{0}\oplus V_{3}\oplus V_{6}\oplus V_{9}\oplus\text{\dots , }\Gamma_\kk\left({\mathcal F}_{1}\right)\simeq V_{1}\oplus V_{4}\oplus V_{7}\oplus V_{10}\oplus\dots.
\notag
\end{equation}

Let us prove that $ \Gamma_\kk\left({\mathcal F}_{0}\right) $ and $ \Gamma_\kk\left({\mathcal F}_{1}\right) $ are simple $ {A} $-modules. Indeed, let $ N $
be a proper submodule of $\Gamma_\kk(\mathcal{F}_0)$. If $ u_{\lambda}\in N $, then $ u_{\lambda+3k}=x^{k} u_{\lambda}\in N $ for all positive $ k $.
Choose the minimal $ \lambda $ such that $ u_{\lambda}\in N $. Then the quotient module has a
decomposition $ V_{\lambda-3}\oplus\dots \oplus V_{0} $, hence it is finite-dimensional. Since $A$ has no non-zero finite-dimensional modules, this is a contradiction. The case of  $\Gamma_\kk(\mathcal{F}_1)$ is very similar. In this way we obtain that, if $ M_{x}={\mathcal F}_{0} $ or $ {\mathcal F}_{1} $, then $ M $ is respectively isomorphic to $ \Gamma_\kk\left({\mathcal F}_{0}\right) $ or $ \Gamma_\kk\left({\mathcal F}_{1}\right) $.

Finally, we show that $ \Gamma_\kk\left({\mathcal F}_{2}\right)=0 $. It is sufficient to check that there is no
non-zero $v\in{\mathcal F}_{2} $ with $F\cdot v=0 $ and
\begin{equation}
H \cdot v=\left(-3k-2\right)v\text{ for }k\in{\mathbb Z}_{\geq0}.
\label{equ23}\end{equation}\myLabel{equ23,}\relax
Indeed, then $v$ would be a solution of the differential equation
\begin{equation}
3y v_{x}=v_{yy}.
\notag\end{equation}
Since $ v\in{\mathcal F}_{2} $,
\[
v = g \left( x,y \right) \exp{} \left( -\frac{y^{3}}{3x} \right)
\]
for some $g\left( x,y \right) \in {\mathbb C} \left[x,x^{-1},y\right]$ such that
\begin{equation}
3y g_{x}=g_{yy}-2\frac{y^{2}}{x}g_{y}-2\frac{y}{x}g.
\notag\end{equation}
As $ g\left(x,y\right) $ is homogeneous with respect to $ H $, we may assume without loss of
generality that
\begin{equation}
g\left(x,y\right)=\sum_{i=0}^{l} b_{i} x^{p-i}y^{3i+s},
\notag\end{equation}
where $ s\in{\mathbb Z}_{\geq0} $, $ p\in{\mathbb Z} $, $ b_{i}\in{\mathbb C} $, $ b_{0}=1 $. The equation on the highest term with respect to $ x $ gives the condition
\begin{equation}
\partial_{y}^{2}\left(y^{s}\right)=0,
\notag\end{equation}
or, equivalently, $ s=0,1$. But $ H\cdot g=\left(3p+s+2\right)g $, hence $ H \cdot v=\left(3p+s+2\right)\cdot v $. Therefore
\begin{equation}
H \cdot v=\left(3p+2\right) v\text{ or }H\cdot v=\left(3p+3\right) v,
\notag\end{equation}
and~\eqref{equ23} does not hold.\end{proof}

\refth{th600} together with \refle{le103} yield the following.

\begin{corollary} \label{cor36}\myLabel{cor36}\relax
In the principal case, up to isomorphism, there are exactly four simple $(\gg,\kk)$-modules with central character $\chi(\threehalfs,\half)$. They have the following $\kk$-module decompositions:
\begin{equation}\label{eq72}
V_{0}\oplus V_{6}\oplus V_{12}\oplus\text{\dots , }V_{1}\oplus V_{7}\oplus V_{13}\oplus\text{\dots , }V_{3}\oplus V_{9}\oplus V_{15}\oplus\text{\dots , }V_{4}\oplus V_{10}\oplus V_{16}\oplus\dots .
\end{equation}
\end{corollary}

\section{$\kk$-characters of simple  bounded ($\sp(4)$, $\sl(2))$-modules}\label{secKchars4-2}
\subsection{The root case.}
In this case, the four simple modules of \refcor{cor35} are nothing but the simple highest weight modules $L_{\threehalfs,\half}$, $L_{\threehalfs,-\half}$, and their respective restricted duals $L'_{\threehalfs,\half}$, $L'_{\threehalfs,-\half}$, i.e. the simple $\bb$-lowest weight modules with lowest weights $(-\threehalfs,-\half)$ and $(-\threehalfs, \half)$. Therefore, by Corollaries \ref{cor101}, \ref{cor102} we conclude that all simple bounded $(\gg,\kk)$-modules are precisely $L_{a,b}$ and the lowest weight modules $L_{-a,-b}'$, where $a>|b|\in\half+\ZZ$. Since $c(L_{a,b})=c(L_{-a,-b}')$, it suffices to compute $c(L_{a,b})$, for $a,b$ as above.

The $\hh$-character of $L_{a,b}$ is given by the formula
\begin{equation}\label{eq61}
\ch_\hh L_{a,b}=\frac{(x^{a-b}-x^{b-a})(y^{a+b}-y^{-a-b})}{(x-x^{-1})(y-y^{-1})(xy-x^{-1}y^{-1})(x^{-1}y-xy^{-1})},
\end{equation}
where $x=e^{\frac{\eps_1-\eps_2}{2}}$, $y=e^{\frac{\eps_1+\eps_2}{2}}$. We rewrite \refeq{eq61} as
\begin{equation}\label{eq92}
\frac{(x^{a-b}-x^{b-a})(y^{a-b}-y^{b-a})}{(x-x^{-1})(y-y^{-1})}y^{-2}(1-x^2y^{-2})^{-1}(1-x^{-2}y^{-2})^{-1}.
\end{equation}

Next we note that
\begin{equation}\label{eq93}
(1-x^2y^{-2})^{-1}(1-x^{-2}y^{-2})^{-1}=\sum_{k=0}^{\infty}y^{-2k}(x^{2k}+x^{2k-4}+\dots+x^{-2k}),
\end{equation}
and use the expression
\[
z^k=x^k+x^{k-2}+\dots+x^{-k}=\frac{x^{k+1}-x^{-(k+1)}}{x-x^{-1}}
\]
to rewrite the right-hand side of \refeq{eq93} in the form
\[
\sum_{k=0}^{\infty}y^{-2k}(z^{2k}-z^{2k-2}+\dots+(-1)^k)=\frac{1}{1+y^{2}}\sum_{k=0}^\infty z^{2k}y^{-2k}.
\]
Now \refeq{eq92} becomes
\[
\ch_\hh L_{a,b}=z^{a-b-1}\frac{y^{a+b}-y^{-a-b}}{y-y^{-1}}\frac{1}{1+y^{2}}\sum_{k=0}^\infty z^{2k}y^{-2k}.
\]
To find the $\kk$-character of $L_{a,b}$, we set $y=1$:
\begin{equation}\label{eq94}
c(L_{a,b})=\frac{a+b}{2}z^{a-b-1}\otimes\sum_{k=0}^\infty z^{2k}.
\end{equation}
Thus, equation \refeq{eq94} implies the following result.
\begin{theorem}~
\begin{itemize}
\item[(a)] If $a-b$ is even and $a+b$ is odd, then
\[
c(L_{a,b})=\frac{a+b}{2}(2z+4z^3+\dots+(a-b)z^{a-b-1}+(a-b)z^{a-b+1}+\dots).
\]
\item[(b)] If $a-b$ is odd and $a+b$ is even, then
\[
c(L_{a,b})=\frac{a+b}{2}(1+3z^2+5z^4+\dots+(a-b)z^{a-b-1}+(a-b)z^{a-b+1}+\dots).
\]
\item[(c)] In the case (a) the minimal $\kk$-type is $V_1$ and its multiplicity is $a+b$. In the case (b) the minimal $\kk$-type is $V_0$ and its multiplicity is $\frac{a+b}{2}$.
\item[(d)] For sufficiently large $i$,
\[
c_i(L_{a,b})=c_{i+2}(L_{a+b})=\frac{(a^2+b^2)(1+(-1)^{a+b-i})}{4}.
\]
\item[(e)] $L_{a,b}$ is $\kk$-multiplicity-free if and only if $a=\threehalfs$, hence the only simple multiplicity-free $(\gg,\kk)$-modules are those with central character $\chi(\threehalfs, \half)$, i.e. the four $\gg$-modules from \refcor{cor36}.
\end{itemize}
\end{theorem}

\subsection{The principal case.}
We now proceed to calculating the $\kk$-characters of all simple bounded $(\gg,\kk)$-modules where $\gg=\sp(4)$ and $\kk$ is the principal subalgebra of $\gg$ fixed in \refse{se6}. In this case, let $M_{\threehalfs,\half}^0$ and $M_{\threehalfs,\half}^1$ denote the simple bounded $(\gg,\kk)$-modules with central character $\chi (\threehalfs,\half)$ and respective $\kk$-module decompositions $V_0\oplus V_6\oplus V_{12}\oplus\dots$ and  $V_1\oplus V_7\oplus V_{13}\oplus\dots$. We set $M_{a,b}^s \eqdef T_{a\eps_1+b\eps_2}^{\threehalfs\eps_1+\half\eps_2}(M_{\threehalfs,\half}^s)$ for $a,b\in\half+\ZZ, a>|b|$, $s\in\{0,1\}$, and $M_{a,b}^s:=0$ for $a,b\in\half+\ZZ, a\leq|b|$, $s\in\{0,1\}$. By $V_{p,q}$ we denote the simple finite-dimensional $\gg=\sp(4)$-module with $\bb$-highest weight $p\eps_1+q\eps_2$ ($p,q\in \ZZ_{\geq 0}$, $p\geq q$).
\begin {lemma}\label{le200}
We have
\begin{equation}\label{eq1155}
V_{1,0}\otimes M_{a,b}^s \simeq M_{a+1,b}^s\oplus M_{a,b+1}^s\oplus M_{a-1,b}^s\oplus M_{a,b-1}^s,
\end{equation}
and, for $a\neq |b|+1$,
\begin{equation}\label{starstar}
V_{1,1}\otimes M_{a,b}^s \simeq M_{a+1,b+1}^s\oplus M_{a,b}^s\oplus M_{a-1,b+1}^s\oplus M_{a+1,b-1}^s\oplus M_{a-1,b-1}^s.
\end{equation}
If $a=b+1,b>0$, then
\begin{equation}\label{star1star}
V_{1,1}\otimes M_{a,b}^s\simeq M_{a+1,b+1}^s\oplus M_{a+1,b-1}^s
\oplus M_{a-1,b-1}^s,
\end{equation}
and if $a=-b+1,b<0$, then
\begin{equation}\label{star2star}
 V_{1,1}\otimes M_{a,b}^s\simeq M_{a+1,b+1}^s\oplus M_{a+1,b-1}^s
\oplus M_{a-1,b+1}^s.
\end{equation}
\end{lemma}
\newcommand{\mabs}{\MM_{a,b}^s}
\begin{proof}
Let us first prove \refeq{eq1155}. In what follows we use the notation of [PS2], Section 5. Let
$\MM_{a,b}^s\eqdef\DD_{G/B}^{a,|b|}\otimes_{\univ^{\chi(a,b)}}M_{a,b}^s$
be the localization of $M_{a,b}$ on $G/B$.
Then as a sheaf of $\univ$-modules $V_{1,0}\otimes\MM_{a,b}^s$ has a filtration of length 4 with the following associated factors given in increasing order:
\[
\OO(-\eps_1)\otimes_\OO \mabs, \quad \OO(-\eps_2)\otimes_\OO\mabs,\quad
\OO(\eps_2)\otimes_\OO\mabs,\quad \OO(\eps_1)\otimes_\OO\mabs.
\]
Note that $Z_\univ$ acts via a character on any of the four associated
 factors, and that these characters are pairwise distinct.
 Therefore, as a sheaf of $\univ$-modules, $V_{1,0}\otimes\mabs$ is isomorphic to the direct sum
\[
\left(\OO(-\eps_1)\otimes_\OO \mabs\right) \oplus \left(\OO(-\eps_2)
\otimes_\OO\mabs\right)\oplus\left( \OO(\eps_2)\otimes_\OO\mabs\right)
\oplus(\OO(\eps_1)\otimes\mabs).
\]
Now we calculate $\Gamma(G/B,V_{1,0}\otimes\mabs)$. If $a=b+1,b>0$, then
$$\Gamma(G/B,\OO(-\eps_1)\otimes_\OO \mabs)=
\Gamma(G/B,\OO(\eps_2)\otimes_\OO \mabs)=0$$
 as there are no bounded modules with these central characters. Similarly, if
$a=-b+1$, $b<0$, then
$$\Gamma(G/B,\OO(-\eps_1)\otimes_\OO \mabs)=
\Gamma(G/B,\OO(-\eps_2)\otimes_\OO \mabs)=0.$$
In all other cases
$$\Gamma(G/B,\OO(\pm\eps_1)\otimes_\OO \mabs)\simeq M^s_{a\pm 1,b},$$
$$\Gamma(G/B,\OO(\pm\eps_2)\otimes_\OO \mabs)\simeq M^s_{a,b \pm 1}.$$
Thus, \refeq{eq1155} is established.


Consider \refeq{starstar}. Then as a sheaf of $\univ$-modules $V_{1,1}\otimes\MM_{a,b}^s$ has a filtration of length 5 with the following associated factors given in increasing order:
\[
\OO(-\eps_1-\eps_2)\otimes_\OO \mabs, \quad \OO(\eps_1-\eps_2)\otimes_\OO\mabs, \quad \mabs,
\]
\[
\OO(-\eps_1+\eps_2)\otimes_\OO\mabs,\quad \OO(\eps_1+\eps_2)\otimes_\OO\mabs.
\]
Note that $Z_\univ$ acts via a character on any of the five associated
factors, and that these characters are pairwise distinct if $a\neq
|b|+1$. Therefore the proof of \refeq{starstar} is very similar to that of \refeq{eq1155}.

Let now $a=b+1$. Then $\mabs$ and $\OO(-\eps_1+\eps_2)\otimes_\OO\mabs$
both afford the central character $\chi (a,b)$. Thus, as a sheaf of
$\univ$-modules,
$V_{1,1}\otimes\mabs$ is isomorphic to the direct sum
\begin{equation}\label{eqle112}
\left(\OO(-\eps_1-\eps_2)\otimes_\OO \mabs\right) \oplus \left(\OO(\eps_1-\eps_2)\otimes_\OO\mabs\right)\oplus \left(\mabs\right)' \oplus
\end{equation}
\[
\oplus\left( \OO(\eps_1+\eps_2)\otimes_\OO\mabs\right),
\]
where for $\left(\mabs\right)'$ we have an exact sequence
\[
0\to\mabs\to\left(\mabs\right)'\to\OO(-\eps_1+\eps_2)\otimes_\OO\mabs\to 0.
\]
We will show that $\Gamma(G/B,(\mabs)')=0$. It suffices to show that
the tensor product $V_{1,1}\otimes M^s_{a,b}$ has no simple constituent with central character $\chi (a,b)$. Indeed, from \refeq{eq1155}, we see that $V_{1,0}\otimes V_{1,0}\otimes M^s_{a,b}$ has exactly two simple constituents affording the central character $\chi (a,b)$ and that both these constituents are isomorphic to $M^s_{a,b}$. Recall that
\begin{equation*}%\label{eqle112}
V_{1,0}\otimes V_{1,0}\cong V_{2,0}\oplus V_{1,1}\oplus V_{0,0}.
\end{equation*}
Clearly, $V_{0,0}\otimes M^s_{a,b}=M^s_{a,b}$. Furthermore, $V_{2,0}$ is the adjoint representation and therefore the very $\gg$-module structure on $M^s_{a,b}$ defines a non-trivial intertwining operator $V_{2,0}\otimes M^s_{a,b}\to M^s_{a,b}$. Thus, $V_{2,0}\otimes M^s_{a,b}$ must have a constituent isomorphic to $M^s_{a,b}$ and consequently $V_{1,1}\otimes M^s_{a,b}$ has no simple constituent affording the central character $\chi(a,b)$. By taking the global sections of the direct sum \refeq{eqle112} we obtain \refeq{star1star}. The case $a=-b+1$, which leads to \refeq{star2star}, is similar.
\end{proof}

\begin{lemma}\label{le201}
There is the following $\kk$-module decomposition
\begin{equation} \label {e1}
M_{\threehalfs,-\half}^{s}\simeq V_{3+s}\oplus V_{9+s}\oplus V_{15+s}\oplus\dots\text{ .}
\end{equation}
\end{lemma}
\begin {proof}
By \refeq{eq1155},
\[
M_{\threehalfs,\half}^{0}\otimes V_{1,0}\simeq M_{\frac{5}{2},\half}^0\oplus M_{\threehalfs,-\half}^0.
\]
As a $\kk$-module, $V_{1,0}$ is isomorphic to $V_3$. Hence $M_{\threehalfs,\half}^{0}\otimes V_{1,0}$ has a $\kk$-module decomposition
\[
\label{e3}
2V_3\oplus V_5\oplus \dots\quad .
\]
Since $\chi(\threehalfs, -\half)=\chi(\threehalfs,\half)$, $M_{\threehalfs,-\half}^0$ must have one of the four $\kk$-module decompositions \refeq{eq72}, and hence \refeq{eq1155} implies \refeq{e1} for $s=0$. Similarly, $M_{\threehalfs,\half}^1\otimes V_{1,0}$ has the $\kk$-module decomposition $V_2\oplus 2V_4\oplus \dots$, which implies \refeq{e1} for $s=1$.
\end{proof}

We set now $\phi_{a,b}^s(z)\eqdef c(M_{a,b}^s)$ for $a,b\in\half+\ZZ, a\geq|b|$, $s\in\{0,1\}$ and extend the definition of $\phi _{a,b}^s(z)$ to arbitrary pairs $a,b\in \half +\ZZ$ by putting
\begin{equation}\label{eqDelta}
\phi_{a,b}^s(z)=-\phi_{b,a}^s(z)=-\phi_{-b,-a}^s(z)=\phi_{-a,-b}^s(z).
\end{equation}

\begin {lemma}\label{le204}
For all $a,b\in \half+\ZZ$ and $s\in\{0,1\}$,
\[
\pi(\phi_{a,b}^s(z^3+z+z^{-1}+z^{-3}))=\phi_{a-1,b}^s+\phi_{a+1,b}^s+\phi_{a,b+1}^s+\phi_{a,b-1}^s
\]
\[
\pi(\phi_{a,b}^s(z^4+z^2+1+z^{-2}+z^{-4}))=\phi_{a+1,b+1}^s+\phi_{a-1,b+1}^s+\phi_{a+1,b-1}^s+\phi_{a-1,b-1}^s+\phi_{a,b}^s.
\]
(the projection $\pi$ is introduced in \refsec{secRank2Case}).
\end{lemma}
\begin{proof}
Both equalities are straightforward corollaries of \refle{le200} and \refle{le202} (b) if one takes into account the isomorphisms of $\kk$-modules $V_{1,0}\simeq V_3$ and $V_{1,1}\simeq V_4$.
\end{proof}

We define now $\psi_{a,b}^s(z)\in \CC((z))$ via the conditions:
\begin {itemize}
\item[(c1)]$\displaystyle\psi_{a,b}^s(z)(z^3+z+z^{-1}+z^{-3})=\psi_{a+1,b}^s(z)+\psi_{a-1,b}^s(z)+\psi_{a,b+1}^s(z)+\psi_{a,b-1}^s(z)$,
\item[(c2)]$\displaystyle\psi_{a,b}^s(z)(z^4+z^2+1+z^{-2}+z^{-4})=\psi_{a+1,b+1}^s(z)+\psi_{a+1,b-1}^s(z)+\psi_{a-1,b+1}^s(z)+\psi_{a-1,b-1}^s(z)+\psi_{a,b}^s(z)$,
\item[(c3)]$\displaystyle\psi_{a,b}^s(z)=-\psi_{b,a}^s(z)=-\psi_{-b,-a}^s(z)=\psi_{-a,-b}^s(z)$,
\item[(c4)]$\displaystyle\psi_{\threehalfs,\half}^s(z)=\frac{z^s}{1-z^6}, \quad \psi_{\threehalfs,-\half}^{s}(z)=\frac{z^{3+s}}{1-z^6}$.
\end{itemize}

\begin {theorem}\label{th206}
The Laurent series $\psi_{a,b}^s(z)$ exists and is unique. Moreover,
\begin{equation}\label{eq4}
\psi_{a,b}^{s}(z)=\frac{ z^{5+s}(z^{3a+b}-z^{a+3b}-z^{-a-3b}+z^{-3a-b})-z^{6+s}(z^{3a-b}-z^{-a+3b}-z^{a-3b}+z^{-3a+b})}{(1-z^2)^2(1-z^4)(1-z^6)}.
\end{equation}
\end{theorem}
\begin{proof}
We show first that $\psi_{a,b}^s(z)$ is unique if it exists. By \refeq{eqDelta} $\psi_{a,b}^s(z)$ is determined by $\psi_{a,b}^s(z)$ for $a>|b|$. Assume, by induction on $a$, that $\psi_{a,b}^s(z)$  is unique for all $a\leq a_0$, $|b|<a$. Then equation (c1) determines $\psi_{a_0+1,b}^s(z)$, and equation (c2) determines $\psi_{a_0+1,a_0}^s(z)$ and $\psi_{a_0+1,a_0+1}^s(z)$.

To prove the existence of $\psi_{a,b}^s(z)$, it suffices to verify that the right-hand side of \refeq{eq4} satisfies all conditions (c1)-(c4). This is a direct calculation, which is simplified by the observation that both Laurent polynomials
\[
z^{3a+b}-z^{a+3b}-z^{-a-3b}+z^{-3a-b},
\]
\[
z^{3a-b}-z^{-a+3b}-z^{a-3b}+z^{-3a+b}
\]
satisfy (c1),(c2) and (c3). The condition (c4) is satisfied only by the entire expression.
\end{proof}

\begin {corollary}\label{cor207}
\[
\phi_{a,b}^s=\pi(\psi_{a,b}^s).
\]
\end{corollary}

\begin {corollary}\label{cor207'}
Any simple bounded \gkm~ is either even or odd. More precisely, $M_{a,b}^s$ is even if $a+b+s$ is even, and $M_{a,b}^s$ is odd if $a+b+s$ is odd.
\end{corollary}

In the calculations below we use binomial coefficients $\binom{s}{k}$,
for which we always assume $\binom{s}{k}=0$ if $s$ or $k$ are not integers.

\begin {lemma}\label{le208}
\[
\frac{1}{(1-z^2)^2(1-z^4)(1-z^6)}=\sum_{n=0}^{\infty}\gamma(n)z^{2n},
\]
where
\[
\gamma(n)\eqdef \frac{1}{144}\left[ 119\binom{n+3}{3}-179\binom{n+2}{3}+109\binom{n+1}{3}-25\binom{n}{3}\right]+\frac{(-1)^n}{16}+\frac{\beta(n)}{9}
\]
and
\[
\beta(n)\eqdef\triplebrace{0}{n\equiv 1~(\mod 3)}{1}{n\equiv 0~(\mod 3)}{-1}{n\equiv -1~(\mod 3)}.
\]
\end{lemma}
\begin{proof}
The statement follows from the identity
\[
\frac{1}{(1-z^2)^2(1-z^4)(1-z^6)}=\frac{119-179z^2+109z^4-25z^6}{144(1-z^2)^4}+\frac{1}{16(1+z^2)}+\frac{1+z^2}{9(1+z^2+z^4)}.
\]
\end{proof}

\begin {corollary}\label{cor208}
Let
\begin{eqnarray*}
\delta_{a,b}^s(n)&=&\gamma\left(\frac{n-(3a+b+5)-s}{2}\right)-\gamma\left(\frac{n-(a+3b+5)-s}{2}\right)-\\
&&-\gamma\left(\frac{n-(-a-3b+5)-s}{2}\right)+\gamma\left(\frac{n-(-3a-b+5)-s}{2}\right)-\\
&&-\gamma\left(\frac{n-(3a-b+6)-s}{2}\right)+\gamma\left(\frac{n-(-a+3b+6)-s}{2}\right)+\\
&&+\gamma\left(\frac{n-(a-3b+6)-s}{2}\right)-\gamma\left(\frac{n-(-3a+b+6)-s}{2}\right).
\end{eqnarray*}
Then
\[
c_i(M_{a,b}^s)=\delta_{a,b}^s(i)-\delta_{a,b}^s(-i-2).
\]
\end{corollary}
\begin{proof}
The statement follows directly from \refth{th206}, \refcor{cor207}, and \refle{le208}.
\end{proof}

\begin{corollary}\label{cor209}
For any simple bounded \gkm~ $M$, $c_i(M)=c_{i+6}(M)$ for sufficiently large $i\in\NN$.
\end{corollary}
\begin{proof}
The given \gkm~ $M$ is isomorphic to $M_{a,b}^s$ for some $a,b\in \half+\ZZ$, $s\in\{0,1\}$. For sufficiently large $i$, $\delta_{a,b}^s(-i-2)=0$, hence $c_i(M)=\delta_{a,b}^s(i)$. The explicit formula for $\gamma(i)$ from \refle{le208} implies that $\delta_{a,b}^s(i+6n)$ is a polynomial in $n$. Since this polynomial is a bounded function, it is necessarily a constant.
\end{proof}

For large enough values of $i$, \refcor{cor209} enables us to write $c_{\overline{i}}(M_{a,b}^s)$, $\overline{i}\in \ZZ_6$. Here are simple explicit expressions for $c_{\overline{i}}(M_{a,b}^s)$.

\begin {theorem}\label{thX}
Let $\sigma_{a,b}\eqdef\triplebrace{1}{\text{ if }3|2a,3\nmid 2b}{-1}{\text{ if }3|2b,3\nmid 2a}{0}{\text{ in all other cases}}$.

Then
\flushleft{
$\displaystyle
c_{\overline{0+s}}(M_{a,b}^s)=\frac{1}{6}(1+(-1)^{a+b})\left(\frac{a^2-b^2}{2} +2\sigma_{a,b}\right),
$
}
\flushleft{
$\displaystyle
c_{\overline{1+s}}(M_{a,b}^s)=c_{\overline{5+s}}(M_{a,b}^s)=\frac{1}{6}(1-(-1)^{a+b})\left(\frac{a^2-b^2}{2} -\sigma_{a,b}\right),
$
}
\flushleft{
$\displaystyle
c_{\overline{2+s}}(M_{a,b}^s)=c_{\overline{4+s}}(M_{a,b}^0)=\frac{1}{6}(1+(-1)^{a+b})\left(\frac{a^2-b^2}{2} -\sigma_{a,b}\right),
$
}
\flushleft{
$\displaystyle
c_{\overline{3+s}}(M_{a,b}^s)=\frac{1}{6}(1-(-1)^{a+b})\left(\frac{a^2-b^2}{2} +2\sigma_{a,b}\right).
$
}

\end{theorem}
\begin{proof}
Let $\{\xi_{\bar i}\}_{\bar i\in\ZZ_6}$ denote the standard basis in $\CC^6$. Set
\[
\overline{\phi}_{a,b}^s\eqdef\sum_{\overline{i}\in\ZZ_6}c_{\overline{i}}(M_{a,b}^s)\xi_{\bar
  i}
\]
for $a,b\in\half+\ZZ$, $a\geq |b|$. Extend $\overline{\phi}_{a,b}^s$ to all $a,b\in\half+\ZZ$ by putting
\[
\overline{\phi}_{a,b}^s =-\overline{\phi}_{b,a}^s=-\overline{\phi}_{-b,-a}^s=\overline{\phi}_{-a,-b}^s,
\]
and let $S,T:\CC^6\to \CC^6$ be the linear operators
\[
S(\xi_{\overline{i}})\eqdef 2\xi_{\overline{i+3}}+\xi_{\overline{i+1}}+\xi_{\overline{i-1}},\quad T(\xi_{\overline{i}})\eqdef 2\xi_{\overline{i+2}}+2\xi_{\overline{i+4}}.
\]
Then $\overline{\phi}_{a,b}^s$ satisfy the following version of conditions (c1)-(c4):
\begin{itemize}
\item[(c5)] $S(\overline{\phi}_{a,b}^s )=\overline{\phi}_{a+1,b}^s +\overline{\phi}_{a,b+1}^s +\overline{\phi}_{a-1,b}^s +\overline{\phi}_{a,b-1}^s$,
\item[(c6)] $T(\overline{\phi}_{a,b}^s )=\overline{\phi}_{a+1,b+1}^s +\overline{\phi}_{a-1,b+1}^s +\overline{\phi}_{a+1,b-1}^s +\overline{\phi}_{a-1,b-1}^s$,
\item[(c7)] $\overline{\phi}_{a,b}^s =-\overline{\phi}_{b,a}^s =-\overline{\phi}_{-b,-a}^s =\overline{\phi}_{-a,-b}^s $,
\item[(c8)] $\overline{\phi}_{\threehalfs,\half}^s=\xi_{\overline{s}},\quad  \overline{\phi}_{\threehalfs,-\half}^s=\xi_{\overline{3+s}}$.
\end{itemize}
Denote by $\omega$ a primitive sixth root of unity. Then $\displaystyle\{\eta_{\overline{i}}\eqdef\sum_{\overline{j}\in \ZZ_6}\omega^{\overline{i}\overline{j}}\xi_{\overline{j}}\}_{\overline{i}\in\ZZ_6}$ is an eigenbasis for $S$ and $T$.
Put
\[
\eta_{\overline{0},a,b}\eqdef \frac{(a^2-b^2)}{2}\eta_{\overline{0}},\quad \eta_{\overline{3},a,b}\eqdef (-1)^{a+b}\frac{(a^2-b^2)}{2}\eta_{\overline{3}},
\]
\[
\eta_{\overline{2},a,b}\eqdef \sigma_{a,b}\eta_{\overline{2}},\quad \eta_{\overline{4},a,b}\eqdef \sigma_{a,b}\eta_{\overline{4}},
\]
\[
\eta_{\overline{3},a,b}\eqdef (-1)^{a+b}\sigma_{a,b}\eta_{\overline{3}},\quad \eta_{\overline{5},a,b}\eqdef (-1)^{a+b}\sigma_{a,b}\eta_{\overline{5}}.
\]
Using the identity
\[
\sigma_{a,b}=\frac{\omega^{2b}+\omega^{-2b}-\omega^{2a}-\omega^{-2a}}{3},
\]
one can easily check that $\eta_{\overline{i},a,b}$ satisfies (c5)-(c7). The linear combination
\[
\overline{\phi}_{a,b}^s=\frac{1}{6}\sum_{\overline{i}\in\ZZ_6}\omega^{-\overline{is}}\eta_{\overline{i},a,b}
\]
satisfies the condition (c8), hence its coefficients in the basis
$\{\xi_{\bar i}\}$ equal $c_{\overline{i}}\left(M_{a,b}^s\right)$.
\end{proof}

\begin{corollary}\label{cor712}
The following is a complete list of multiplicity-free simple \gkm s: $M_{\threehalfs,\pm\half}^{s}$,  $M_{\frac{5}{2},\pm \threehalfs}^s$, $M_{\frac{5}{2},\pm\half}^{s}$,$M_{\frac{7}{2},\pm\frac{5}{2}}^{s}$, $s\in\{0,1\}$.
\end{corollary}
\begin {proof}
A straightforward computation based on \refth{thX} shows that $c_{\overline{i}}(M_{a,b}^s)\in\{0,1\}$ for $\overline{i}\in\ZZ_6$ iff $(a,b)$ is one of the pairs $\displaystyle\left(\frac{3}{2},\pm\frac{1}{2} \right)$, $\displaystyle\left(\frac{5}{2},\pm\frac{3}{2} \right)$, $\displaystyle\left(\frac{5}{2},\pm\frac{1}{2} \right)$, and $\displaystyle\left(\frac{7}{2},\pm\frac{5}{2} \right)$. Then, using \refcor{cor208} one verifies that all modules $M_{a,b}^s$ for $(a,b)$ as above are indeed multiplicity-free.
\end{proof}

\begin {theorem}\label{th813}~

\begin {itemize}
\item[(a)] The minimal $\kk$-type of any even (respectively, odd) bounded simple \gkm~M equals $V_0$, $V_2$ or $V_4$ (resp., $V_1$ or $V_3$).
\item[(b)] If $M$ is an even (respectively, odd) simple module in $\BB^{\chi(a,b)}$, then $c_0(M)$ (resp., $c_1(M)$) equals $\displaystyle \frac{a\pm b}{6}+\eps$ or
$\displaystyle \frac{a\pm b}{12}+\eps$ (resp., $\displaystyle \frac{a\pm b}{3}+\eps$ or $\displaystyle \frac{a\pm b}{6}+\eps$) for some $\eps$ with $|\eps|<1$.
\end{itemize}
\end{theorem}
\begin{proof}
(a) Note that for any bounded \gkm~M, $c_i(M)$ equals the constant term of the Laurent polynomial $z^{-i}(1-z^{2i+2})c(M)$. Hence $c_1(M)+c_3(M)$ equals the constant term in the Laurent expansion of $(z^{-1}(1-z^4)+z^{-3}(1-z^8))c(M)$. A straightforward calculation shows that for $M=M_{a,b}^s$ the latter is nothing but the constant term of the Laurent series
\[
\frac{z^{3a+b+2+s}-z^{a+3b+2+s}-z^{-a-3b+2+s}+z^{-3a-b+2+s}-z^{-3a+b+3+s}}{(1-z^2)^3}+
\]
\[
+\frac {z^{a-3b+3+s}+z^{-a+3b+3+s}-z^{3a-b+3+s}}{(1-z^2)^3}.
\]
Using the identity
\begin{equation}\label{eq851}
\frac{1}{(1-z^2)^3}=\sum_{n=0}^{\infty}\binom{n+2}{2}z^{2n},
\end{equation}
we obtain
\begin{equation}
\label{eqX}
c_1(M_{a,b}^s)+c_3(M_{a,b}^s)\eqdefrev d_{a,b}^s=\binom{\frac{-3a-b+2-s}{2}}{2}-\binom{\frac{-a-3b+2-s}{2}}{2}-\binom{\frac{a+3b+2-s}{2}}{2}+\binom{\frac{3a+b+2-s}{2}}{2}
\end{equation}
\[
-\binom{\frac{3a-b+1-s}{2}}{2}+\binom{\frac{-a+3b+1-s}{2}}{2}+\binom{\frac{a-3b+1-s}{2}}{2}-\binom{\frac{-3a+b+1-s}{2}}{2},
\]
where we set $\binom{l}{2}:=0$ for $l\notin\ZZ_{\geq 0}$.

This expression is a piecewise polynomial function which equals identically zero whenever $M_{a,b}^s$ is even, i.e. when $a+b+s$ is even. In fact, the right hand side of \refeq{eqX} turns out to be very simple as an explicit calculation shows that, for $a+b+s$ odd,
\begin{equation}\label{eq813}
d_{a,b}^s=\doublebrace{\displaystyle\frac{a+(-1)^{s+1}b}{2}}{\mathrm{ for~} a+(-1)^s3b\geq 0}{a+(-1)^{s}b}{\mathrm{ for~} a+(-1)^s3b\leq 0}.
\end{equation}
Since $a>|b|$, the right hand side of \refeq{eq813} is never 0, i.e. the minimal $\kk$-type of $M_{a,b}^s$ is $V_1$ or $V_3$ whenever $a+b+s$ is odd.

A similar analysis proves that the minimal $\kk$-type of $M_{a,b}^s$ is $V_0$, $V_2$, or $V_4$ whenever $a+b+s$ is even. Indeed, in this case
\[
e_{a,b}^s\eqdef c_0(M_{a,b}^s)+c_2(M_{a,b}^s)+c_4(M_{a,b}^s)
\]
equals the constant term of the Laurent series
\[
(1-z^2)+z^{-2}(1-z^6)+z^{-4}(1-z^{10})c(M)\quad .
\]
Using the identity
\[
\frac{(1-z^2)+z^{-2}(1-z^6)+z^{-4}(1-z^{10})}{(1-z^2)^2(1-z^4)(1-z^6)}=\frac{1}{8z^4}\left( \frac{7+4z^2+z^4}{(1-z^2)^3}+\frac{1}{(1+z^2)}\right),
\]
as well as the identity \refeq{eq851}, we calculate
\begin{eqnarray*}
e_{a,b}^s&=& \theta\left( \frac{-3a-b-1-s}{2} \right)-\theta\left( \frac{-a-3b-1-s}{2} \right)-\\
&&-\theta\left( \frac{a+3b-1-s}{2} \right)+\theta\left( \frac{3a+b-1-s}{2} \right)-\\
&&-\theta\left( \frac{3a-b-2-s}{2} \right)+\theta\left( \frac{-a+3b-2-s}{2} \right)+\\
&&+\theta\left( \frac{a-3b-2-s}{2} \right)-\theta\left( \frac{-3a+b-2-s}{2} \right),
\end{eqnarray*}
where $\displaystyle\theta(n)\eqdef \frac{3}{4}n^2+\frac{3}{2}n+\frac{7}{8}+\frac{(-1)^n}{8}$ for $n\in\ZZ_{\geq 0}$ and $\theta(n):=0$ otherwise. Further calculations show:
\begin{equation}\label{eqeabs}
e_{a,b}^s=\doublebrace{\displaystyle \frac{3}{4}\left(a+(-1)^{s+1}b\right)+\frac{(-1)^ {\frac{a+(-1)^{s+1}b-1}{2}}}{4}} {\mathrm{~for~}(-1)^sa+3b\geq 0}{\displaystyle\frac{3}{2}\left(a+(-1)^sb\right)}{\mathrm{~for~}(-1)^sa+3b\leq 0}
\end{equation}
under the assumption that $a+b+s$ is even. Since the right-hand side of \refeq{eqeabs} never equals 0, we obtain that $e_{a,b}^s\neq 0$ under the same assumption. Hence the minimal $\kk$-type of any even simple bounded \gkm~ equals $V_0$, $V_2$, or $V_4$.

(b) To compute $c_0(M)$ we use the identity
\begin{eqnarray*}
\frac{1-z^2}{(1-z^2)^2(1-z^4)(1-z^6)}&=&\frac{1}{(1-z^2)(1-z^4)(1-z^6)}\\
&=&\frac{47-52z^2+17z^4}{72(1-z^2)^3}+\frac{1}{8(1+z^2)}+\frac{2-z^2-z^4}{9(1-z^6)}
\end{eqnarray*}
which yields
\begin{eqnarray*}
c_0(M_{a,b}^s)&=&\gamma'\left( \frac{-3a-b-5-s}{2}\right)-\gamma'\left( \frac{-a-3b-5-s}{2}\right)-\\
&&-\gamma'\left( \frac{a+3b-5-s}{2}\right)+\gamma'\left( \frac{3a+b-5-s}{2}\right)-\\
&&-\gamma'\left( \frac{3a-b-6-s}{2}\right)+\gamma'\left( \frac{-a+3b-6-s}{2}\right)+\\
&&+\gamma'\left( \frac{a-3b-6-s}{2}\right)-\gamma'\left( \frac{-3a+b-6-s}{2}\right),
\end{eqnarray*}
where
\[
\gamma'(n)\eqdef
\frac{n^2}{12}+\frac{n}{2}+\frac{94}{144}+\frac{(-1)^n}{8}+\frac{\sigma'(n)}{9},
\]
\[
\sigma'(n)\eqdef\doublebrace{-1}{3\nmid n}{2}{3\mid n}
\]
for $n\in \ZZ_{\geq 0}$ and $\gamma'(n)=\sigma'(n):=0$ otherwise. Similarly, using the identity
\begin{eqnarray*}
\frac{z^{-1}(1-z^4)}{(1-z^2)^2(1-z^4)(1-z^6)}=z^{-1}\left( \frac{8-7z^2+2z^4}{9(1-z^2)^3}+\frac{1+z^2-2z^4}{9(1-z^6)}\right)
\end{eqnarray*}
we obtain
\begin{eqnarray*}
c_1(M_{a,b}^s)&=&\gamma''\left( \frac{-3a-b-4-s}{2}\right)-\gamma''\left( \frac{-a-3b-4-s}{2}\right)-\\
&&-\gamma''\left( \frac{a+3b-4-s}{2}\right)+\gamma''\left( \frac{3a+b-4-s}{2}\right)-\\
&&-\gamma''\left( \frac{3a-b-5-s}{2}\right)+\gamma''\left( \frac{-a+3b-5-s}{2}\right)+\\
&&+\gamma''\left( \frac{a-3b-5-s}{2}\right)-\gamma''\left( \frac{-3a+b-5-s}{2}\right),
\end{eqnarray*}
where
\[
\gamma''(n)\eqdef \frac{n^2}{6}+\frac{5n}{6}+\frac{8}{9}+\frac{\sigma''(n)}{9},
\]
\[
\sigma''(n)\eqdef\doublebrace{-2}{n= -1 (\mod 3)}{1}{n\neq -1 (\mod 3)}
\]
for $n\in \ZZ_{\geq 0}$ and $\gamma''(n)=\sigma''(n):=0$ otherwise. Using the expressions for $c_0(M^s_{a,b})$ and $c_1(M^s_{a,b})$ we notice that the terms $\frac{(-1)^n}{8}+\frac{\sigma'(n)}{9}$ and $\frac{\sigma''(n)}{9}$ will give a contribution $\eps$ with $|\eps|<1$. Thus, a direct computation implies
\[
c_0(M^s_{a,b})=\doublebrace{\frac{a+(-1)^sb}{6}+\eps}{\mathrm{~for~}a+(-1)^s3b<0}{\frac{a-(-1)^sb}{12}+\eps}{\mathrm{~for~}a+(-1)^s3b>0},
\]
\[
c_1(M^s_{a,b})=\doublebrace{\frac{a-(-1)^sb}{6}+\eps}{\mathrm{~for~}a+(-1)^s3b>0}{\frac{a+(-1)^sb}{3}+\eps}{\mathrm{~for~}a+(-1)^s3b<0.}
\]
\end{proof}

\begin{corollary}
For $a\pm b\geq 24$, the minimal $\kk$-type of $M_{a,b}^s$ equals $V_0$ (respectively, $V_1$) if $a+b+s$ is odd (resp., even).
\end{corollary}

\begin{corollary}
A simple \gkm~ with minimal $\kk$-type $V_i$ for $i\geq 5$ is unbounded.
\end{corollary}

Note that all simple \gkm s of finite type over $\kk$ with minimal $\kk$-type $V_i$ for $i\geq 6$ are classified in \cite{PZ2}. In particular it is proved, \cite{PZ2}, that if $M$ is a \gkm~ with minimal $\kk$-type $V_i$ for $i\geq 6$, then $M$ is necessarily of finite type over $\kk$ and $c_i(M)=1$. Recently G. Zuckerman and the first named author have shown that this holds also for $i=5$, and \refth{th813} (b) implies that the statement is false for $i\leq 1$.

%%%%%%%%%%%%%%%%%%%%%%%%%%%%%%%%%
%%%%%%%%%%%%%%%%%%%%%%%%%%%%%%%%%
%%%%%%%%%%%%%%%%%%%%%%%%%%%%%%%%%
%%%%%%%%%%%%%%%%%%%%%%%%%%%%%%%%%
%%%%%%%%%%%%%%%%%%%%%%%%%%%%%%%%%
%%%%%%%%%%%%%%%%%%%%%%%%%%%%%%%%%
%%%%%%%%%%%%%%%%%%%%%%%%%%%%%%%%%
%%%%%%%%%%%%%%%%%%%%%%%%%%%%%%%%%
%%%%%%%%%%%%%%%%%%%%%%%%%%%%%%%%%
%%%%%%%%%%%%%%%%%%%%%%%%%%%%%%%%%
%%%%%%%%%%%%%%%%%%%%%%%%%%%%%%%%%
%%%%%%%%%%%%%%%%%%%%%%%%%%%%%%%%%
%%%%%%%%%%%%%%%%%%%%%%%%%%%%%%%%%
%%%%%%%%%%%%%%%%%%%%%%%%%%%%%%%%%
%%%%%%%%%%%%%%%%%%%%%%%%%%%%%%%%%
%%%%%%%%%%%%%%%%%%%%%%%%%%%%%%%%%
%%%%%%%%%%%%%%%%%%%%%%%%%%%%%%%%%
%%%%%%%%%%%%%%%%%%%%%%%%%%%%%%%%%
\begin{thebibliography}{20}
\bibitem[B]{B} V. Bargmann, Irreducible unitary representations of the Lorentz group, Ann. of Math. (2) \textbf{48} (1947), 586--640.
\bibitem[BG]{BG} J. Bernstein, S. Gelfand, Tensor products of finite and infinite-dimensional representations of semi-simple Lie algebras, Compositio Math. \textbf{41} (1980), 245--285.
\bibitem[Dix]{Dix} J. Dixmier, Enveloping algebras, North-Nolland Math. Library, 14, North Holland, Amsterdam, 1977.
\bibitem[F]{F} S. L. Fernando, Lie algebra modules with finite-dimensional weight spaces I, Trans. Amer. Math. Soc. \textbf{322} (1990), 757--781.
\bibitem[Fo]{Fo} A. I. Fomin, Quasisimple irreducible representations of the group $SL(3,\mathbb R)$. Func. Anal. i Pril. \textbf{9} (1975), No 3, 67--74.
\bibitem[GN]{GN} I. M. Gelfand, M. A. Naimark, Unitary representations of the Lorentz group, Izvestiya Akad. Nauk SSSR, Ser. Mat. \textbf{11} (1947), 411--504.
%\bibitem[GP]{GP} I. M. Gelfand, V. A. Ponomarev, The classification of
% indecomposable infinitesimal representations of the Lorentz group,
% Dokl. Acad. Nauk USSR 176 (1967), 502-505.
\bibitem[HC]{HC} Harish-Chandra, Infinite irreducible representations of the Lorentz group, Proc. Roy. Soc. London, Ser. A \textbf{189} (1947), 372--401.
%\bibitem[H]{H} A. van den Hombergh, Harish-Chandra modules and
 % representations of step algebras, Thesis Nijmegen University (1976).
%\bibitem[J]{J} A. Joseph, Some ring theoretic techniques and open
% problems in enveloping algebras, Noncommutative Rings, (Berkeley,
 % Calif., 1989),  MSRI Publ. 24, Springer Verlag, New York, 1992.
%\bibitem[K]{K} M. Kashiwara, B-functions and holonomic systems. Rationality of roots of B-functions. Invent. Math. \textbf{38} (1976/77), 33--53
%\bibitem[K1]{K1} V. G. Kac, Some remarks on nilpotent of orbits, J. Algebra \textbf{64} (1980), 190--213.
%\bibitem[K2]{K2} V. G. Kac, Constructing groups associated to infinite-dimensional Lie algebras, Infinite-dimensional groups with applications (Berkeley, Calif., 1984), MSRI Publ. 4, Springer-Verlag, New York, 1985.
%\bibitem[KL]{KL} G.R. Krause, T.H. Lenagan, Growth of algebras and Gelfand-Kirillov dimension, Pitman, Boston-London-Melbourne, 1985.
\bibitem[Kn1]{Kn1} F. Knop, Der Zentralisator einer Liealgebra in
einer einhuellenden Algebra (German), J. Reine Angew. Math.  \textbf{406} (1990), 5--9.
\bibitem[Kn2]{Kn2} F. Knop, Classification of multiplicity-free symplectic  representations.  J. Algebra  \textbf{301} (2006), 531--553.
\bibitem[KV]{KV} A. W. Knapp, D. Vogan, Jr., Cohomological induction and unitary representations, Princeton University Press, Princeton, NJ, 1995.
%\bibitem[M]{M} O. Mathieu, Classification of irreducible weight
 % modules, Ann. Fourier \textbf{50} (2000),  537--592.
\bibitem[OV]{OV} A. Onishchik, E. Vinberg, Lie groups and algebraic
  groups. Springer series in Soviet mathematics, Springer-Verlag,
  Berlin, 1990.
\bibitem[PS1]{PS1} I. Penkov, V. Serganova, Generalized Harish-Chandra modules, Moscow Math. J. \textbf{2} (2002), 753--767.
\bibitem[PS2]{PS2} I. Penkov, V. Serganova, On bounded generalized Harish-Chandra modules, preprint, 2009, math.jacobs-university.de/penkov
\bibitem[PSZ]{PSZ} I. Penkov, V. Serganova, G. Zuckerman, On the existence of \gkm s of finite type, Duke Math. J. \textbf{125} (2004), 329-349.
\bibitem[PZ1]{PZ1} I. Penkov, G. Zuckerman, Generalized Harish-Chandra modules: a new direction of the structure theory of representations, Acta Applic. Math. \textbf{81} (2004), 311--326.
\bibitem[PZ2]{PZ2} I. Penkov, G. Zuckerman, Generalized Harish-Chandra modules with generic minimal $\kk$-type, Asian J. of Math. \textbf{8} (2004), 795-812.
\bibitem[PZ3]{PZ3} I. Penkov, G. Zuckerman, A construction of generalized Harish-Chandra  modules with arbitrary minimal $\kk$-type, Canad. Math. Bull. 50 (2007), 603-609.
\bibitem[S]{S} Dj. Sijacki, All $SL(3,\mathbb R)$ ladder
  representations, J. Math. Phys. \textbf{31} (1990), 1872--1876.
%\bibitem[Sch]{Sch}W. Schmid, Die Randwerte holomorpher Funktionen auf
%  hermitesch symmetrischen Raeumen (German),
% Invent. Math. \textbf{116}(1969/1970), 61--80.
\bibitem[Sp]{Sp} B. Speh, The unitary dual of $GL(3,\mathbb R)$ and
  $GL(4,\mathbb R)$, Ann. Math.\textbf{258}(1981/82), 113--133.  
%\bibitem[Str]{Str} R.S. Strichartz, Harmonic analysis on  hyperboloids,
%J. Fuct. Analysis \textbf{12}(1973), 341--383.
\bibitem[V1]{V1} D. Vogan, Jr., Singular unitary representations. In Lecture Notes in Mathematics 880, Springer-Verlag, 1981.
\bibitem[V2]{V3} D. Vogan, Jr., The unitary dual of
  $G_2$. Invent. Math. \textbf{116} (1994), 677--791.
%\bibitem[VK]{VK} E.B. Vinberg, B.N. Kimelfield, Homogeneous domains on flag manifolds and spherical subgroups of semi-simple Lie groups, Func. Anal. i Pril. \textbf{12} (1978), 12--19.
%\bibitem[Z]{Z} G. Zuckerman, Tensor product of finite and
% infinite-dimensional representations of semisimple Lie
%  groups, Ann. Math. \textbf{106}(1977), 295--308.


\end{thebibliography}

\end{document}
